\documentclass[11pt,twoside,openany]{book}  % allow chapters to start on either side (no forced blank pages)
% Core Packages
\usepackage[table]{xcolor}
\usepackage{fontspec}
% Font Setup (fonts located in /content/fonts as provided)
\setmainfont{EBGaramond}[
    Path = /content/fonts/EBGaramond/,
    UprightFont = *-SemiBold,
    ItalicFont = *-Italic,
    BoldFont = *-Bold,
    BoldItalicFont = *-BoldItalic,
    Scale = 1.05
]

\setmonofont{RobotoMono}[
    Path = /content/fonts/RobotoMono/static/,
    UprightFont = *-Regular,
    ItalicFont = *-Italic,
    BoldFont = *-Bold,
    BoldItalicFont = *-BoldItalic,
    Scale = 0.88
]
\usepackage{colortbl}
\usepackage[letterpaper, twoside, twocolumn,
    inner=0.75in, outer=0.5in,
    top=0.5in, bottom=1.0in,
    columnsep=0.35in]{geometry}
\addtolength{\textheight}{-1\baselineskip}
\usepackage{longtable}
\usepackage{tabularx}
\usepackage{array}
\usepackage{booktabs}
\usepackage{microtype}
\usepackage{parskip}
\usepackage{setspace}
\usepackage{xparse}
\usepackage{eso-pic}
\usepackage{tcolorbox}
\tcbuselibrary{skins}
\usepackage{enumitem}
\usepackage{wrapfig}
\usepackage{ifthen}
\usepackage{atbegshi}
\usepackage{float}
\usepackage{needspace}

% Paragraph Settings
\setlength{\parindent}{1em}
\setlength{\parskip}{0pt}
\clubpenalty=10000
\widowpenalty=10000


% Color Definitions (DndRed and others for styling)
\definecolor{DndRed}{RGB}{153, 0, 0}
\definecolor{Parchment}{RGB}{253, 246, 227}
\definecolor{StatblockBack}{RGB}{200, 150, 90}
\definecolor{sectiontitle}{RGB}{139, 0, 0}
\definecolor{TableHeader}{RGB}{184,160,122}
\definecolor{AltRow}{RGB}{220, 237, 204} % Light sage green
\definecolor{TableRowAlt}{RGB}{191, 145, 60} % brass
\definecolor{CalloutBoxBack}{RGB}{220, 237, 204}
\definecolor{CalloutBoxBorder}{RGB}{100, 160, 80}

% Use fancyhdr for custom headers/footers
\usepackage{fancyhdr}
\setlength{\headheight}{14pt}
\pagestyle{fancy}
	\fancyhf{}
	\fancyhead[LE]{\scshape\itshape\color{DndRed}{\leftmark}}
	\fancyhead[RO]{\scshape\itshape\color{DndRed}{\rightmark}}
	\fancyfoot[LE,RO]{\color{DndRed}{\thepage}}
	\renewcommand{\headrulewidth}{0pt}
	\renewcommand{\footrulewidth}{0pt}

\fancypagestyle{plain}{%
  \fancyhf{}
  \fancyhead[LE]{\scshape\itshape\color{DndRed}{\leftmark}}
  \fancyhead[RO]{\scshape\itshape\color{DndRed}{\rightmark}}
  \fancyfoot[LE,RO]{\color{DndRed}{\thepage}}
  \renewcommand{\headrulewidth}{0pt}
  \renewcommand{\footrulewidth}{0pt}
}

% List Customization (unchanged)
\newlist{fiveitemize}{itemize}{2}
\setlist[fiveitemize,1]{
    label=\textbullet,
    left=0pt,
    labelsep=0.5em,
    itemsep=2pt,
    parsep=0pt,
    topsep=4pt,
    partopsep=0pt,
    listparindent=0pt,
    font=\normalsize
}
\setlist[fiveitemize,2]{
    label=–,
    left=1.5em,
    labelsep=0.5em,
    itemsep=2pt,
    parsep=0pt,
    topsep=2pt,
    partopsep=0pt,
    listparindent=0pt,
    font=\normalsize\itshape
}

% Tcolorbox Definitions (unchanged)
\newtcolorbox{StatblockBackground}{
    enhanced,
    colback=StatblockBack,
    colframe=black!0,
    arc=2mm,
    boxrule=0pt,
    boxsep=4pt,
	left=6pt, right=6pt, top=4pt, bottom=4pt,
    sharp corners=south,
    width=\columnwidth,
    halign=center,
    valign=center,
	fontupper=\small,
}
\newtcolorbox{CalloutBoxSub}[1]{
    enhanced,
    colback=CalloutBoxBack,
    colframe=CalloutBoxBorder,
    arc=2mm,
    boxrule=1pt,
    boxsep=6pt,
    fontupper=\normalsize\itshape,
    width=\linewidth,
    before upper={\centering\textbf{\Large\color{DndRed}#1}\par\vspace{0.5em}},
	after=\vspace{-\baselineskip},
}
\newtcolorbox{CalloutBox}[1]{
    enhanced,
    colback=CalloutBoxBack,
    colframe=CalloutBoxBorder,
    arc=2mm,
    boxrule=1pt,
    boxsep=6pt,
    fontupper=\normalsize\itshape,
    width=\linewidth,
    before upper={\centering\textbf{\Large\color{DndRed}#1}\par\vspace{0.5em}},
}
\newtcolorbox{PurchasingTableBox}{
    enhanced,
    colback=white,
    colframe=black!0,
    boxrule=0pt,
    boxsep=0pt,
    width=0.85\textwidth,
    left=8pt, right=8pt, top=8pt, bottom=8pt,
}
\newtcolorbox{ClassProgressionBox}{
	enhanced,
	colback=white,
	colframe=black!0,
	boxrule=0pt,
	boxsep=0pt,
	width=0.85\textwidth
	left=8pt, right=8pt, top=8pt, bottom=8pt,
}
\newenvironment{PurchasingTable}{
    \begin{PurchasingTableBox}
    \begin{longtable}{
        >{\centering\arraybackslash}p{2.15cm}
        >{\centering\arraybackslash}p{2.15cm}
        >{\centering\arraybackslash}p{1.8cm}
        >{\centering\arraybackslash}p{1.4cm}
        >{\centering\arraybackslash}p{1.4cm}
        >{\centering\arraybackslash}p{1.4cm}
        >{\centering\arraybackslash}p{1.4cm}
        >{\centering\arraybackslash}p{1.8cm}
    }
}{
    \end{longtable}
    \end{PurchasingTableBox}
}
\newtcolorbox{SmallTableBox}{
  colback=Parchment,
  colframe=black,
  boxrule=0.5pt,
  arc=1pt,
  boxsep=0pt,
  left=6pt,
  right=6pt,
  top=4pt,
  bottom=4pt,
  enhanced,
}
\newtcolorbox{FancyTableBox}{
  enhanced,
  frame hidden,
  interior style={left color=Parchment, right color=Parchment!90!black},
  overlay={
    \draw[line width=0.5pt, black] (frame.north west) rectangle (frame.south east);
  },
  boxrule=0.5pt,
  arc=1pt,
  width=\linewidth,
}


\usepackage{titlesec}
% Part title format: centered, Huge, bold, red. Numbered parts include "Part <Roman>:"
\titleformat{\part}[display]
{\centering\Huge\bfseries\color{DndRed}}
{\partname~\thepart}{1em}{}
\titlespacing*{\part}{0pt}{3\baselineskip}{2em}
% Chapter title format: hang, Large, bold, small caps, red.
% Main matter will show "Chapter N: Title".
\titleformat{\chapter}[hang]
{\Huge\bfseries\scshape\color{DndRed}}
{\chaptername~\thechapter}{1em}{}
\titlespacing*{\chapter}{0pt}{2\baselineskip}{1em}
% Section title format: Large, bold, red. Numbering in Roman numerals (main matter).
\renewcommand{\thesection}{\Roman{section}}
\titleformat{\section}
{\LARGE\bfseries\color{DndRed}}
{\thesection}{0.5em}{}
\titlespacing*{\section}{0pt}{1.5\baselineskip}{0.5em}
\renewcommand{\thesubsection}{\thesection.\arabic{subsection}}
\titleformat{\subsection}
{\Large\bfseries\itshape\color{DndRed}}
{\thesubsection}{0.5em}{}
\titlespacing*{\subsection}{0pt}{1.2\baselineskip}{0.5em}
\titleformat{\subsubsection}
{\normalsize\bfseries\color{DndRed}}
{\thesubsubsection}{0.5em}{}
\titlespacing*{\subsubsection}{0pt}{0.8\baselineskip}{0.5em}
  
% --------------------------------------------------
% PACKAGES REQUIRED FOR FULL-WIDTH IMAGE + PATCHES
% --------------------------------------------------
\usepackage{etoolbox}
\usepackage{tikz}
\usepackage{transparent}
\usepackage{varwidth}
\usetikzlibrary{calc,fadings}
\pgfdeclareverticalshading{topfade}{1cm}{%
  color(0cm)=(white);          % fully faded at very top
  color(0.9cm)=(black) % fully visible by 1cm down
}
\pgfdeclarefading{imageTopFade}{%
  \pgfuseshading{topfade}
}
\usepackage{graphicx}
\usepackage{balance} % optional: balances column bottom height
\usepackage{multicol}
\setlength{\columnsep}{0.35in}
\setlength{\columnseprule}{0pt}


% --------------------------------------------------
% FULL-WIDTH IMAGE AT TOP OF PAGE
% Usage: \FullWidthImageTop[<height>]{<filename>}
% --------------------------------------------------
\newcommand{\ChapterSplashPage}[3][0.4\textheight]{%
  \clearpage
  \onecolumn
  \thispagestyle{plain}% allow fancyhdr styling without headers on top of image
  \refstepcounter{chapter}%
  \addcontentsline{toc}{chapter}{#3}%
  \markboth{Chapter \thechapter: #3}{}%

  \begin{tikzpicture}[remember picture, overlay]
    \node[anchor=north west, inner sep=0pt] at (current page.north west)
      {\includegraphics[width=\paperwidth,height=#1]{#2}};%
  \end{tikzpicture}%

  \vspace*{#1}
  {\Huge\bfseries\scshape\color{DndRed} Chapter \thechapter: #3\par}
  \vspace{1.5em}
}
\newcommand{\SectionSplashPage}[2][0.4\textheight]{%
  \clearpage
  \onecolumn
  \thispagestyle{plain}% allow fancyhdr styling without headers on top of image
  \begin{tikzpicture}[remember picture, overlay]
    \node[anchor=north west, inner sep=0pt] at (current page.north west)
      {\includegraphics[width=\paperwidth,height=#1]{#2}};%
  \end{tikzpicture}%

  \vspace*{#1}
  \vspace{1.5em}
}


% --------------------------------------------------
% OPTIONAL: CLEANER TOC DOT SPACING
% --------------------------------------------------
\makeatletter
\patchcmd{\@starttoc}{\@dotsep}{\@ldotsep}{}{}
\makeatother



\usepackage{hyperref}
% Hyperref Setup (bookmarks open and numbered)
\hypersetup{
    colorlinks=true,
    linkcolor=red,
    urlcolor=blue,
    pdfstartview=FitH,
    bookmarksopen=true,
    bookmarksnumbered=true
}

\usepackage{tocloft}
% Ensure part and chapter numbering in TOC
\renewcommand{\cftpartpresnum}{}            % no "Part " label in TOC
\renewcommand{\cftpartaftersnum}{\quad}     % space after the Roman numeral
\renewcommand{\cftchappresnum}{Chapter~}    % prefix "Chapter " for chapters in TOC
\renewcommand{\cftchapaftersnum}{:}         % colon after chapter number

% Adjust TOC indent/spacing for numbered labels (widths for alignment)
\settowidth{\cftpartnumwidth}{\bfseries\color{DndRed}IV\quad}    % width for part num (e.g., "IV")
\settowidth{\cftchapnumwidth}{\bfseries\color{DndRed}Chapter 10:} % width for chapter num (e.g., "Chapter 10:")

% Use dot leaders for chapters, sections, subsections
\renewcommand{\cftchapleader}{\cftdotfill{\cftdotsep}}
\renewcommand{\cftsecleader}{\cftdotfill{\cftdotsep}}
\renewcommand{\cftsubsecleader}{\cftdotfill{\cftdotsep}}

% Set TOC entry fonts and colors
\renewcommand{\cftpartfont}{\bfseries\scshape\color{DndRed}}
\renewcommand{\cftchapfont}{\bfseries\scshape\color{DndRed}}
\renewcommand{\cftsecfont}{\bfseries\itshape\color{DndRed}}
\renewcommand{\cftsubsecfont}{\itshape\color{DndRed}}
% (Optional) also color the page numbers in TOC
\renewcommand{\cftpartpagefont}{\color{DndRed}}
\renewcommand{\cftchappagefont}{\color{DndRed}}
\renewcommand{\cftsecpagefont}{\color{DndRed}}
\renewcommand{\cftsubsecpagefont}{\color{DndRed}}
% Wrapfigure already loaded earlier if needed for other images
\newsavebox{\StormforgedBox}
\newlength{\StormforgedWidth}
\newsavebox{\WarsmithBox}
\newlength{\WarsmithWidth}



% Booktabs Red Line Commands (unchanged)
\newcommand{\RedToprule}{\arrayrulecolor{DndRed}\specialrule{1.5pt}{0pt}{2pt}\arrayrulecolor{black}}
\newcommand{\RedThickMidrule}{\arrayrulecolor{DndRed}\specialrule{1.2pt}{2pt}{2pt}\arrayrulecolor{black}}
\newcommand{\RedThinMidrule}{\arrayrulecolor{DndRed}\specialrule{0.8pt}{2pt}{2pt}\arrayrulecolor{black}}
\newcommand{\RedBottomrule}{\arrayrulecolor{DndRed}\specialrule{1.5pt}{2pt}{0pt}\arrayrulecolor{black}}
\newcommand{\MonsterGradientRule}{
  \noindent
  \begin{tikzpicture}
    \fill[fill=DndRed, draw=none] (0,0) -- ++(\linewidth,0) -- ++(0,-0.05) -- (0,-0.2) -- cycle;
  \end{tikzpicture}
}
\newcommand{\MonsterGradientThin}{
  \noindent
  \begin{tikzpicture}
    \fill[fill=DndRed, draw=none] (0,0) -- ++(\linewidth,0) -- ++(0,-0.05) -- ++(-\linewidth,0) -- cycle;
  \end{tikzpicture}
}
\newcommand{\ThinRule}{
  \noindent
  \begin{tikzpicture}
    \fill[fill=DndRed, draw=none] (0,0) -- ++(\linewidth,0) -- ++(0,-0.05) -- ++(-\linewidth,0) -- cycle;
  \end{tikzpicture}
}
\newcommand{\Subheading}[1]{%
  \vspace{0.8\baselineskip}%
  {\noindent\color{DndRed}\scshape #1\par}%
  \vspace{0.5em}%
  \ThinRule%
  \vspace{1pt}%
}
% Monster Macros (unchanged)

\newcommand{\StatEntry}[2]{%
  \multicolumn{1}{>{\raggedright\arraybackslash}X}{{\color{DndRed}\textbf{#1}}} &
  \multicolumn{5}{X}{#2} \\[1.5pt]}
\newcommand{\StatLine}[6]{{\color{DndRed}\textbf{#1}} & 
{\color{DndRed}\textbf{#2}} & {\color{DndRed}\textbf{#3}} & 
{\color{DndRed}\textbf{#4}} & {\color{DndRed}\textbf{#5}} & 
{\color{DndRed}\textbf{#6}} \\}
\newcommand{\MonsterSection}[1]{%
  \vspace{4pt}
  \MonsterGradientThin
  {\raggedright\color{DndRed}\normalsize\bfseries\scshape #1}
  \par\vspace{2pt}
  \MonsterGradientThin
}
\newcommand{\MonsterTrait}[2]{%
  \noindent\raggedright\textbf{#1.} #2\par\vspace{4pt}
}
\newcommand{\MonsterAction}[2]{%
  \noindent\raggedright\textbf{#1.} #2\par\vspace{4pt}
}

% Small and Full-Width Image Commands (unchanged)
\newcommand{\SmallWrappedImage}[3][0.4\textwidth]{
  \begin{wrapfigure}{r}{#1}
    \vspace{-10pt}
    \includegraphics[width=#1]{#2}
    \caption*{\small\textit{#3}}
  \end{wrapfigure}
}

% Global page background (75% opacity parchment image on every page)
\AddToShipoutPicture{%
    \begin{tikzpicture}[remember picture,overlay]
        \node[opacity=0.75,inner sep=0pt] at (current page.center) {\includegraphics[width=\paperwidth,height=\paperheight]{Bakku_Background.png}};
    \end{tikzpicture}%
}

% Full-page image command (background not cleared; image overlays on page)
\newcommand{\FullPageImageNoBackground}[1]{%
  \onecolumn
  \begingroup
    \thispagestyle{empty}%
    \noindent
    \begin{tikzpicture}[remember picture,overlay]
      \node[inner sep=0pt] at (current page.center)
        {\includegraphics[width=\paperwidth,height=\paperheight]{#1}};%
    \end{tikzpicture}%
    \null
    \vfill
    \break % Force page end but avoid LaTeX’s float balancing
  \endgroup
  \twocolumn
}
\newcommand{\FullPageImageNoBG}[1]{%
  \begin{titlepage}
    \thispagestyle{empty}
    \begin{tikzpicture}[remember picture,overlay]
      \node[inner sep=0pt] at (current page.center)
        {\includegraphics[width=\paperwidth,height=\paperheight]{#1}};
    \end{tikzpicture}
    \null
  \end{titlepage}
}


\begin{document}
\pagenumbering{gobble}

% **Cover Page** – full-bleed cover image, no headers/footers
\FullPageImageNoBG{Malkoros_Crafting_Cover_Page.png}
\FullPageImageNoBackground{Malkoros_Map.png}

\frontmatter % Roman numerals for front matter, does NOT reset colors
\pagenumbering{roman}

% Preface
\part*{Preface}


\chapter*{Foreword}
\addcontentsline{toc}{chapter}{Foreword}
\markboth{Foreword}{}


\section*{Introduction}
\addcontentsline{toc}{section}{Introduction}
\markright{Introduction}
Hail, Adventurer(s)! If you are reading this, that means that you have found Bakku's Crafting Mechanics and Other Creations, and I thank you for your patronage, whether it be by coin or piracy, I care not and welcome your use of this information. I have put a great deal of work into this supplement for 5th Edition Dungeons and Dragons, and my most sincere hope is that it grants insight, widsom, or endless hours of fun for you!

This book is a log of the thousands upon thousands of tests, templates, and trials that I have played or tested. It is the log of the thousands of characters you have played, and the countless experiences you have shared with Dungeons and Dragons that helped me craft the contents within. I hope for only one thing with this book: that you, holding this book, also get to have the same fun that I had along the way.

I have loved working on this book. Not because I love fiddling with layouts or parsing hundreds of thousands of lines of code (though, I really do enjoy it), and not even because of my love for the game system it is for, but because I knew that at the end of the process I could share it with a community of folks that would enjoy it. I love that I get to make something people love to use. Hearing the stories of your characters, the journeys they went on, and the stories that were brought to life is why I do this. It is why this book exists, and it is why that however and wherever you got this book, I hope you get to use it to have fun. Please share your experiences with me via Fantasy Grounds forum!
\section*{Thank You}
\addcontentsline{toc}{section}{Thank You}
On that note, use this book however it works for you. It is your book. Do you love every feature as it is written? That’s fantastic. Do you want to tweak it for your games? Go ahead! That’s what this is all about! I saw a game I loved, but also something else I wanted to do with it. If you see something you love here, do what you want with it. There is no wrong way to use this book, as long as you and your group are having fun.

There are three individuals that I would like to send special thanks to, and those would be my wife, and my two very good friends who play with me in person. I have dedicated a ridiculous amount of my personal time to this supplement, and my wife stood by me and supported me through every day. She is the most selfless person in this world through my eyes and has unwavering support for me in everything that I do; I only hope those who are reading this are so lucky to find a similar individual in their own lives. My two best friends (I don't get out much, okay...don't judge, Adventurer) serve as further inspiration for creating this. Though we can read millions of comments from people around the world who enjoy Dungeons and Dragons, it is the experience of actually playing with those like-minded people that unlocks the magic within this game, and I would not have realized my love and passion for this craft, had I not had a consistent table to play at with them. A very special thank you to Jessie, Tyler, and Ryne... I love you, guys (and girl)! Finally, to the millions that have played and enjoyed Dungeons and Dragons. To all of you... thank you for being part of this.

I don’t know what folks usually write in these forewords. I probably should have looked that up before I wrote this. But I’ve said what I wanted to say, so... stop reading this little blurb of rambling text, and get out there and have fun playing the game!

-Bakku
\vfill
\null
\clearpage
\FullPageImageNoBackground{Intro.png}
\clearpage

 
\chapter*{Content}
\addcontentsline{toc}{chapter}{Content}
\markboth{Content}{}
The character options in this book are vast. The number of classes does not fully convey the vista of opportunity presented. As I write this summary, I've come fresh off the heels of seeing more than forty characters created with these two classes, and playing up to level 6, and can say that no two of the characters were truly alike.

This is the wonder and potential of modular content. This introduction to the character options is meant to be a digestible preview that will help you find where you want to dive into these new options.
\section*{Chapter 1: Classes}
\addcontentsline{toc}{section}{Chapter 1: Classes}
\markright{Chapter 1: Classes}
There are two classes introduced in this supplement as of the date you are reading this. 
\begin{fiveitemize}
	\item The Inventor
	\item The Psion
\end{fiveitemize}
\subsection*{Inventor}
\addcontentsline{toc}{subsection}{Inventor}
The inventor is a new class that concerns itself with the creation of custom loadouts of magic items tailored to your character. It is the adventuring version of a magical craftsman, combining expert tool use, magic, and creativity to unleash a new way to play.

The inventor offers in-depth and modular customization in a way few other class options do, and is most suited to veteran players, but is fundamentally designed to play alongside the default options in a balanced manner—the depth of the class comes from its character creation options, not from its turn-to-turn gameplay.

The inventor is a class that is defined by their subclasses and choices. A modular toolbox that lets you build your own class, and puts a great deal of care into making sure those potential outcomes are as balanced as possible. Here’s a quick rundown of what to expect from each subclass:
\begin{fiveitemize}
	\item Gadgetsmith. A classic gadgeteer, these are highly adaptable characters with impressive mobility. They excel at utility and disrupting their foes. They serve a role closest to a rogue. A gadgetsmith always has a trick up their sleeve, and is almost as good at getting out of a tight spot with a smoke bomb as they are at getting into trouble in the first place. Consider this option if you thrive on the frustration of your foes, and always want some options at your fingertips, but don’t care about maximizing your damage and durability.
	\item Golemsmith. Where the golem is the star of the show. Take to the field as a smart bloke with a powerful minion that does (most) of your fighting for you, evolving to work in tandem as a tag team duo of danger, or specialize in making the most powerful golem you can make. You have vast opportunities to customize your golem, bringing your vision to life for a balanced combination of durability and damage. Consider this option if you view the art of perfecting a golem buddy as what you want your character to be about.
	\item Infusionsmith. An enchanter that brings their magic to life, choosing from three different routes: a classic wandslinger with infused magical wands that freely tosses out low-level spells and cantrips with abandon, an animator with a retinue of floating weapons, or an infused swordsman, serving as a weapon enhancing magical hybrid, generally serving as a utility reach striker. Consider this option if you want to unleash damage on your foes with a bevy of magical armaments.
	\item Potionsmith. An alchemist who delves into the secrets of reactions and infusions. Wielding a wide range of options from poisoned bolts to explosive reactions, they serve as a mid range striker or support... or a bit of both. Consider this option if you have an unfortunate fascination with explosives, like to prepare for the challenges ahead, or want to provide a copious amount of healing to your allies.
	\item Thundersmith. The cannoneer or weaponsmith, these are the ones that specialize in forging the biggest, baddest weapons. From gunslingers to wielders of massive kinetic hammers, they devote their art form to a single perfect creation. Consider this option if you want to play a striker of any range.
	\item Warsmith. An armor-focused inventor who builds themselves into a man of iron... either in the metaphorical sense of encasing themselves in a magical suit of armor, or... more directly. Highly specialized options, and capable of great feats of strength. Consider this option if you want to take the role of a moderately durable frontline striker, or an adaptable intelligence character.
	\item Fleshsmith. The doctor of fleshwarping... or the fleshwarped monstrosity they have created. Endless horrifying options to delve the depths of creativity turned inward. Pick this option if you want an unsettling support or durable front line brawler.
	\item Cursesmith. A dealer of powers best left undealt. The dark reflection of Infusionsmiths. these wield great but double edged powers. Consider this option if you want to play a striker who is almost as dangerous to their foes as they are to themselves.
	\item Runesmith. Articulates the power of runes, taking the paths of Runic Knights, a durable front line, Runic Mystics, a front line striker, or Runic Sages, a backline support who shares their runic powers.
	\item Relicsmith. Delve into the mystique of the divine and harness the holy. Bring light to darkness and set your foes ablaze with divine fire. Consider this option if you want a divine support striker hybrid, or mid to close range strikers with a holy flare.
\end{fiveitemize}

\subsection*{Psion}
\addcontentsline{toc}{subsection}{Psion}
The psion is a new class that, unsurprisingly, brings psionics to the game. Psionics are a new branch of supernatural power with a long legacy in the game, and have often been featured as one of the core sources of power; just as you might have power that is divine, arcane, or primal in origin, you might have power that is psionic in origin.

Just as with the branches of magic, it can have overlap in functionality, but brings its own focus and powers to the table. The psion is an eminently customizable and specialized sort of character. It can be nearly anything, be it ranged striker, controller, support, or hybrid of various kinds. It has a long standing tradition of excelling at melee-caster roles, and can even acquire the durability to serve as a bulkier frontline character with enough specialization.

The subclasses of the psion focus on which branch of psionics you specialize in and include a suite of abilities that help you master it, but are from the destiny of your character. You can select a second discipline from any in the class, meaning that almost any subclass can achieve almost any role through that combination, but some lend themselves more easily to certain paths. Here’s a quick rundown of what to expect from each subclass:
\begin{fiveitemize}
	\item Awakened Mind. Starting with the power of Telepathy, this option lends itself to being a balanced mid-line character providing control, damage, and utility, making great use of cycling their powers. Pairs great with Telekinesis for maximum control or consumption for higher damage.
	\item Shaper’s Mind. Able to bring their imagination into reality with projection, they offer great damage and situational utility out of the box. Capable of offering some of the highest damage in the class, it can also offer some potent control options when specialized down that route. A bit more complicated in its interactions than some options.
	\item Transcended Mind. An option that easily builds itself into a support, or a self-supporting hybrid style character. They build up a pool of extra damage when buffing their allies (or themselves) which they can expend to deal damage, as well as get a stable of potent buffs from their default branch of psionics, Enhancement. Pairs well with Telekinesis or Psychokinesis for a balanced support striker hybrid, with the freedom to use nearly any discipline as a melee, self-buffing style character.
	\item Unleashed Mind. Starting with the discipline of Telekinesis, this subclass offers a conduit for unleashing the raw power of the mind, with its ever-increasing rampage. Ideal as a ranged striker, it has a flexible array of options for it’s second discipline.
	\item Wandering Mind. Uniquely scaling off Dexterity, this subclass excels at a hybrid striker role. Starting with the discipline of Transposition, they deftly hop about the battlefield, leaving behind rifts in space as they strike with their weapon. As they often focus on Dexterity, they prefer to complement their powers with a discipline that isn’t as dependent on Intelligence, such as Precognition or Enhancement.
	\item Elemental Mind. A psion that is perhaps a little too in touch with the elements, it offers a wide variety of blasting options with its default discipline of Psychokinetics, but is able to specialize in a wide variety of options: a focused ice psion with enhancement makes a surprisingly durable character!
	\item Consuming Mind. That which goes bump in the night, these start with the discipline of consumption, and lend themselves to being strikers. Pairs best with Telepathy, but can work with any discipline that offers them a bit of control to lock down the targets they want to eat.
\end{fiveitemize}

\subsection*{Class Toolbox}
\addcontentsline{toc}{subsection}{Class Toolbox}

Open up the classes and tinker with their guts! Experience the joys of modularity and homebrewing yourself. The toolbox is here to empower GMs and players to make exactly what they want. This toolbox provides even more additional options for both Inventors and Psions as examples of how the classes can be expanded in your hands. If you find yourself wanting even more, the toolbox is the place to delve.

\begin{fiveitemize}
	\item Class Toolbox
\end{fiveitemize}

\section*{Chapter 2: Subclasses}
\addcontentsline{toc}{section}{Chapter 2: Subclasses}
In addition to the classes presented, there’s a selection of subclasses for the default classes that tie into the mechanics of invention and psionics that are added by this book. Each of these subclasses offers a unique spin and blend of some elements of those themes, some more directly than others. Here is a brief summary of what they provide:

\begin{fiveitemize}
	\item Path of the Raging Mind (Barbarian). A psionic option that allows you to exercise that classic barbarian rage to fling things with your mind.
	\item Path of the Exosuit (Barbarian). An option that provides an armored barbarian... of a sort. Bonded to a powerful exosuit, you gain a host of new defenses and strengths with your artificial aid.
	\item College of Thunder (Bard). Leveraging new inventions to bring greater volume, this option gives the bard a bit more damage and a lot more sound.
	\item Mystery Cult (Cleric). Delving into the unknown, this cleric mingles their faith with more mysterious psionic influences, giving them broad (if somewhat sinister) utility.
	\item Circle of Nightmares (Druid). Unlocking the secrets from beyond, this druid is not restricted to drawing inspiration from the material plane for wild shape, they unlock horrifying powers of transformation.
	\item Tech Knight (Fighter). Focusing on the science of war as much as the art of war, this fighter uses dangerous new weapons and gadgets to be the state of the art in combat.
	\item Way of the Soulblade (Monk). Focusing all of their psionic power into a blade, these monks gain the ability to strike with concentrated psionic powers, as well as dip into the broader pool of psionic abilities, freely mixing ki and psionics.
	\item Oath of Sanity (Paladin). Few can face horrors beyond and come out with their sanity intact. These are the first and last line of defense to the madness out there.
	\item Specialist (Ranger). Specialized in enhancing their ammunition, these are a versatile option that always has a trick up their quiver.
	\item Mind Reaper (Ranger). A terrifying reaper that leverages psionic powers to infuse their strikes and infiltrate the minds of their foes.
	\item Gadgeteer (Rogue). A tricky rogue with many tricks, with a handful of new toys that mesh nicely with the rogue kit for a variety of new (mostly explosive) options.
	\item Aether Heart (Sorcerer). A sorcerer whose very heart is suffused with raw arcane power, this arcane sorcerer option specializes in wielding metamagic.
	\item Planetouched (Sorcerer). Wielding the powers of psionic translocation, these sorcerers live at the very edge of reality.
	\item The Ancient Intelligence (Warlock). Making a pact with an intelligence beyond mortal comprehension, these warlocks wield sufficiently advanced technology.
	\item That Which Is Beyond (Warlock). For when old ones don’t quite seem great enough, these warlocks unlock telepathic powers in exchange for their harrowing experiences with that which is beyond.
	\item Order of Creation (Wizard). A wizard that has turned their inventive magic into building useful tools, and an ever present mechanical familiar.
	\item Order of the Unknown (Wizard). For some, magic isn’t enough—out there are the unknown powers, and these wizards won’t rest until they’ve delved them, mastering psionics on top of their magic... with the only price paid being their sanity.
\end{fiveitemize}

\section*{Chapter 3: More Stuff}
\addcontentsline{toc}{section}{Chapter 3: More Stuff}
\markright{Chapter 3: More Stuff}
There is a variety of other character content within these pages such as:
\begin{fiveitemize}
	\item Backgrounds
		\begin{fiveitemize}
			\item Apothecary. A mixer of potions, a master of the subtle art of producing potent magical effects from the right distillations
			\item Engineer. A person with an intellectual bent that can put their abilities to more practical uses: designing everything from seige equipment to bridges.
			\item Tinker. A travelilng vagabond that falls somewhere between a merchant and a sage.
		\end{fiveitemize}
	\item Races
		\begin{fiveitemize}
			\item Ironwrought. The origins of Ironwrought are many and varied. They can be the artisanal masterpiece of a Golemsmith, or perhaps a relic of a mass production of an ancient lost civilization... or perhaps not-so-lost a civilization. Their source and purpose will vary greatly. Were they made with sapience in mind, or did they gain it later? Were they made for labor or war?
			\item Farling Farlings. are a race that comes from beyond the material plane, though few are certain where. While humanoid in shape, they are immediately distinguished from mortal races. They have a slender build, large eyes, and tendrils in place of hair. Their strange nature and inherent psionic abilities make them easily feared, but they tend to be far more friendly to mortal races than most of that which comes from the beyond.
		\end{fiveitemize}
	\item Templates
		\begin{fiveitemize}
			\item Augmented. The augmented are creatures that have replaced parts of the body with mechanical and magical parts—some sort of an arcane cyborg. The origin of these parts is many and varied, as is their nature: anything from living magic to mechanical means is possible, based on the setting and character.
			\item Warped. The warped are members of races that have been deeply influenced by the powers from beyond, twisting their mind and body with mysterious power that has left them fundamentally changed.
		\end{fiveitemize}
	\item Feats
		\begin{fiveitemize}
			\item Feats. General Feats
		\end{fiveitemize}
	\item Class Feats
		\begin{fiveitemize}
			\item Class Feats. Inventor and Psion Feats
		\end{fiveitemize}
	\item Racial Feats
		\begin{fiveitemize}
			\item Racial Feats. Ironwrought and Farling Feats
		\end{fiveitemize}
\end{fiveitemize}

\section*{Chapter 4: Spells}
\addcontentsline{toc}{section}{Chapter 4: Spells}
\begin{fiveitemize}
	\item{Spells}
\end{fiveitemize}

\section*{Chapter 5: Monsters}
\addcontentsline{toc}{section}{Chapter 5: Monsters}
\begin{fiveitemize}
	\item Monster Section
		\begin{fiveitemize}
			\item Aberrant Abstraction
			\item Adaptive Slime
			\item Ancient Sentinel
			\item Clockwork Drake
			\item Eldritch Spawn
			\item Ironwrought
			\item Lurking Maw
			\item Rampaging Imprint
			\item Silver Relics
			\item The Uncaged
			\item Unearthed Destroyer
		\end{fiveitemize}	
	\item NPCs
		\begin{fiveitemize}
			\item Animator
			\item Ascendant Master
			\item Cryokinetic
			\item Doctor
			\item Elemental Adept
			\item Flesh Abomination
			\item Gadgeteer
			\item Golem Mechanic
			\item Gunslinger
			\item Ironwrought Golem
			\item Pyrokinetic
			\item Shaper
			\item Telepath
			\item Thunderer
			\item Warsmith
		\end{fiveitemize}
\end{fiveitemize}

\chapter*{Balance}
\addcontentsline{toc}{chapter}{Balance}
\markboth{Balance}{}
The content of this book is something I’ve worked on extensively to make as balanced as possible, but I think it’s worth talking here in this foreword about what that means, why it might or might not matter, and what to do if your experiences differ.

The point of balance in a cooperative game is twofold. First, so that the GM can run the game within the parameters of the base content without having to tailor and rebalance the experience. Second, so that players can feel free to make the choices they want to make about their character without feeling they are placing themselves at a disadvantage.

The content in this book is different from the default content of the game, but a great deal of care and testing has gone into making sure that it can play alongside the default content. Everything between these covers has been used in games, and reviewed by many eyes... but that guarantees little.

Content can be used in many ways. Expectations can differ. Rolled stats, crazy luck, creative uses, and clever combinations of abilities can always result in varied experiences... particularly when there are as many options as are presented here.

My intention with this book is that the content will feel at home among the default options of the base game and its first expansion. That is what I use in my game in addition to my own content, and that’s what I’ve compared the options presented here to throughout the process of creating it.

\begin{CalloutBox}{Bakku's Note for Players}
When using character options, creativity is encouraged; not only is creativity in the title, but everything contained has many options. This book holds a vast amount of options that can be arranged and rearranged.

But you are not playing against your GM. At the end of the day, the GM has the final say on how any ability works, and what is an option. Your GM has allowed this book, and ultimately anything in it is up to them to allow.

The way to have the most fun with these options is to work collaboratively with what you want to do with your GM on the character you want to make, and how the presented options here will play into that.

The more consideration you put into making sure that your character will fit with the GM’s expectations, the more you can count on them feeling good about allowing you to use the full breadth and depth of content—not just this book, but the wide expanse of options out there. The more they have to double check and vet everything to ensure that all the players have fun at the table, the less they will be able to allow it to be used in the game—quite possibly even content that would otherwise be perfectly balanced.

Balance is a collaboration between players and their GM. Look for combinations that are cool. Look for ways to use the rules that are fun. Even look for ways to make your character feel powerful and strong! But if you think you see something that would make your character broken, talk to your GM. You’re going to need to rely on them to allow it anyway, and both of you are going to be happier knowing if it will be allowed before you decide to play it.
\end{CalloutBox}

\begin{CalloutBox}{Bakku's Note for GMs}
At the end of the day, you have all the power when it comes to balancing the game. I’ve attempted to make this content as easy to allow as possible, but if you find something that doesn’t fit in your game, you know your game better than I do.

I would offer some things to consider, the foremost of which is that 5e is not a perfectly balanced game. I have put in a great deal of effort to make sure that nothing contained in the book eclipses the default options, but at the same time, the default options are a very wide spectrum of content in terms of balance.

If you find an option that seems like it is too much, it might be that you need to adjust it, but it also might be a confluence of events that will solve themselves. If a player is doing more damage than other members of the party, will they still be in a few levels? Have they invested more heavily in doing damage?

Certain kinds of druids find great power early in the game. Certain feats allow for a lot of power. This is in no way a suggestion that nothing should be changed (I change many things within my own game) just a reminder to consider everything in context and the potential ramifications down the road of making changes.
\end{CalloutBox}


\section*{Crafting and Balance}
\addcontentsline{toc}{section}{Crafting and Balance}
\markright{Crafting and Balance}
One of the biggest considerations for balance exists not in the player options at all, but in the crafting options. While the crafting system aims to fairly tightly replicate the traditional rate of acquisition using the presented rules, it is undeniable that the synergy between character options (both default options and ones presented here) and magic items that the players can select exists.

The magic items and their present recipes are presented as tools to make playing the game more fun. It is a tool for both the player and GM, but like any tool, consideration should be placed into how it is used.

Not all magic items exist in all settings. The rules present as wide a selection as possible to lower the burden on the GM, but they always have the ability to define which items can or cannot be made (or need to be tweaked) on a case by case basis.


\section*{Closing Thoughts}
\addcontentsline{toc}{section}{Closing Thoughts}
I do not run everything in this book as written every time. I run what works for my game. I have written things in this book in my attempt to capture an experience as close as possible to the default assumptions of the game, but those are nothing more or less than assumptions—a best fit line for as many games as possible.

As much as possible feel free to stand on these rules and build on them as a foundation, but never let what is written here be a shackle. My only goal in writing this book was to provide a set of tools that will be useful to you in running a game that you and your player enjoy. If any part of these rules contradict that goal for you, purge them with fire.

Lastly, there is a large community of folks that use this content in some degree. While they cannot be found in the pages of this book, feel free to reach out and seek opinions, thoughts, and help with anything presented within!


\clearpage
\onecolumn

\phantomsection
\addcontentsline{toc}{chapter}{Table of Contents}
\markboth{}{}
\renewcommand{\contentsname}{} % suppress default "Contents" title

% Manual Red Title (on same page as ToC)
\vspace*{1.5em}
{\centering
  {\color{DndRed}\scshape\Huge Table of Contents\par}
  \vspace{2em}
}

\tableofcontents

\twocolumn





\mainmatter
\pagenumbering{arabic}


\part{Character Options}


\ChapterSplashPage[0.4\textheight]{Inventor.png}{Classes}
\begin{multicols}{2}
\section{Inventor}
A gnome sits hunched over a workbench in a room cluttered with every sort of tool, carefully drawing the final lines to an intricate rune. With a crackling hum of arcane energy, the completed rune flares with power, and she watches with a smile of pride as the golem comes to life and stands.

An ogre looms large over a dwarf, saliva sloping down its chin as it considers its hearty meal. Its growling anticipation turns into horrified roars of agony as the dwarf unloads terrifying blasts of thunderous power from a small, metal tube.

With a mighty bloom and a crackle of energy, a lone figure lands in the middle of the battlefield. As the smoke clears, a towering presence stands amidst the war-torn ruins clad in glowing, mechanical plate. The luminous visor swivels to inspect the wreckage. It raises a gauntlet, and lightning surges forth.

With a beckoning gesture, the elf summons a swirl of swords leaping from their sheaths and springing to life. The shocked bandits begin to edge away, breaking into a run as the swords exuberantly give chase.

Makers of magic-infused objects, inventors are defined by their innovative nature. Like wizards, they see magic as a complex system waiting to be decoded and controlled through a combination of thorough study and investigation. Inventors, though, focus on creating marvelous, new magical objects. Spells are often too ephemeral and temporary for their tastes. Instead, they seek to craft durable, useful items.
\subsection{Cunning Creators}
Every inventor is defined by a specific craft. Inventors see mastering the basics of a craft as the first step towards true progress, the invention of new methods, and innovative approaches. Inventors are arcane engineers, students of invention and warfare who craft deadly firearms, ingenious gadgets, magical armor, and mechanical beings that they can augment with magic.

All inventors are united by their curious and creative nature. To an inventor, magic is an evolving art with a leading edge of discovery and mastery that pushes further ahead with each passing year.

Inventors value novelty and innovation. This penchant pushes them to seek a life of adventure.

\end{multicols}
\newpage
\twocolumn
A hidden ruin might hold a forgotten magic item or a beautifully crafted mirror perfect for magical enhancement. Inventors win respect and renown among their kind by uncovering new lore or inventing new methods of creation.

\subsection{Boundless Imagination}

An inventor is not the sort to be limited by what already exists or by conventional wisdom. They combine all the tools at their disposal to go beyond—to discover new boundaries and promptly exceed them. Few things can stand in the path of an inventor with a mind to change the world.

Seeking to push beyond all and striving for the new is the unifying call of inventors, but how jealously they guard their secrets can vary greatly. Some push their knowledge upon the worlds, while others hoard it.

\subsection{Rivals and Challenges}

Inventors often acquire rivals—other inventors racing to be the first to unlock the secrets of the world, scrambling to be the one after whom a new magical theory or scientific breakthrough will forever be named. Theft and acquisition of stolen ideas can often become an issue when too many inventors clash over a field of expertise.

Further, inventors quickly make enemies of those who would just as soon see the world unchanged.

They are often viewed as dangerous upstarts or radical troublemakers with their newfangled ways of doing things.

These challenges rarely daunt an inventor for long and most often simply fuel their continued drive to rise above. Inventors are rarely set in their ways, and the hated rival of last week may be a research associate the next.

\subsection{Creating an Inventor}

When creating an inventor character, think about your character’s background and drive for adventure. Does the character have a rival? What drove your character down the path of becoming an inventor? Is it about curiosity and innovation, or is it about the power their inventions bring? Did your character learn from another inventor, or did they receive a vision or flash of inspiration to learn their craft?

Consider how your character interacts with the world and what they represent. Consult with your GM regarding guilds or societies your character might belong to.

\subsubsection{Quick Build}

You can make an inventor quickly by following these suggestions. For the Gadgetsmith, Thundersmith, Cursesmith, or Relicsmith specializations, make your highest ability score Dexterity, followed by Intelligence; for the Golemsmith, Infusionsmith, Potionsmith, or Runesmith specializations, make Intelligence your highest, followed by Constitution or Dexterity; and for Warsmith or Fleshsmith specializations, choose either Strength or Intelligence as your highest, with Constitution as your second-highest ability score. Second, choose a background that emphasizes your connection to a guild or society.
\subsection{Class Features}

As an Inventor you gain the following class features.

\subsubsection{Hit Points}

Hit Dice: 1d8 per inventor level \\
Hit Points at 1st Level: 8 + your Constitution modifier \\
Hit Points at Higher Levels: 1d8 (or 5) + your Constitution modifier per inventor level after 1st

\subsubsection{Proficiencies}

Armor: Light armor, medium armor \\
Weapons: Simple weapons, hand crossbows, heavy crossbows \\
Tools: Thieves’ tools, one other tool of your choice \\
Saving Throws: Constitution, Intelligence \\
Skills: Choose three from Arcana, Deception, History, Investigation, Medicine, Nature, Religion, Sleight of Hand

\subsubsection{Equipment}

You start with the following equipment, in addition to the equipment granted by your background:

\begin{fiveitemize}
  \item (a) a light crossbow and case of 20 bolts or (b) any two simple weapons
  \item (a) scale mail, (b) leather armor, or (c) chain mail (if proficient)
  \item Thieves’ tools and a dungeoneer’s pack
\end{fiveitemize}

If you forgo this starting equipment, as well as the items offered by your background, you start with 4d4 × 10 gp to buy your equipment.

\subsubsection{Inventor Specialization}

At 1st level, you focus your craft on a particular specialization: Cursesmith, Fleshsmith, Gadgetsmith, Golemsmith, Infusionsmith, Potionsmith, Relicsmith, Thundersmith, or Warsmith, each of which are detailed at the end of the class description. Your choice grants you features at 1st level and again at 3rd, 5th, and 14th level.

\subsubsection{Magic Item Analysis}

Starting at 1st level, your understanding of magic items allows you to analyze and understand their secrets. You know the detect magic and identify spells, and you can cast them as a ritual, requiring no material components.

\subsubsection{Tool Expertise}

Starting at 2nd level, your proficiency bonus is doubled for any ability check you make that uses any of the tool proficiencies you gained from this class.

\subsubsection{Arcane Retrofit}

Additionally at 2nd level, you gain the ability to retrofit magical gear. During a long rest you can perform a ritual using any +1, +2, or +3 weapon—excluding artifacts and sentient weapons— to transfer the magic into an inventor weapon (a weapon created by a subclass feature or upgrade). For example, you can turn a +1 longsword and your Impact Gauntlet into a +1 Impact Gauntlet . This includes weapon-like Upgrades that make attack and damage rolls, such as a Warsmith’s Force Blast.

You can’t transfer any properties from a magical weapon besides a bonus to attack and damage rolls, and the original weapon is destroyed in the process.

Additionally, you can convert a set of armor with a magical bonus to AC to a lighter armor type.
\begin{CalloutBoxSub}{Additional Retrofits}
At the GM’s option, this feature can be expanded to do similar tasks. For example, you may be allowed to turn a +1 heavy crossbow into a +1 light crossbow, or the GM may choose to allow this feature to work with other magical armor or weapons, such as converting adamantine chain mail into an adamantine chain shirt.
\end{CalloutBoxSub}
\subsubsection{Spellcasting}
At 2nd level, as part of your study of magic you gain the ability to cast spells. The spells you learn are limited in scope, primarily concerned with modifying creatures and objects or creating items.

\textbf{Spell Slots.} The Inventor table shows how many spell slots you have to cast your inventor spells of 1st level and higher. To cast one of these spells, you must expend a slot of the spell’s level or higher. You regain all expended spell slots when you finish a long rest.

For example, if you know the 1st-level spell longstrider and have a 1st-level and a 2nd-level spell slot available, you can cast longstrider using either slot. \\

\textbf{Spells Known of 1st Level and Higher.} You know three 1st-level spells of your choice from the inventor spell list. 

The Spells Known column of the Inventor table shows when you learn more inventor spells of your choice. Each of these spells must be of a level for which you have spell slots. For instance, when you reach 5th level in this class, you can learn one new spell of 1st or 2nd level.

\begin{CalloutBox}{Inventor's Spellcasting}
An inventor is fundamentally someone that understands and regularly interacts with magic, leading to the ability to apply this knowledge as a spellcaster, but how you want to approach your inventor’s spellcasting is up to you. Feel free to explore other approaches in coordination with your GM. From a functional standpoint, the only requirement is that you have a defined material component or focus for a spell that requires it.
Consider for some spells perhaps using your artisan’s tools as your focus, or perhaps a specific item you’ve made is your material component for a spell. Perhaps instead of a “component pouch” that simply contains all the material components you could need for your spells, you have an “inventor’s tool belt” that contains the various spellcasting assistance contraptions you’ve made to cast your spells. For mechanical purposes, the only thing that matters is that it functions in the same way and is used consistently.
\end{CalloutBox}
Additionally, when you gain a level in this class, you can choose one of the inventor spells you know and replace it with another spell from the inventor spell list, which also must be of a level for which you have spell slots.

\textbf{Spell save DC} = 8 + your proficiency bonus + your Intelligence modifier

\textbf{Spell attack modifier} = your proficiency bonus + your Intelligence modifier

\textbf{Spellcasting Focus.} You can use an arcane focus as a spellcasting focus for your inventor spells.

\newpage
\onecolumn
\noindent
\begin{minipage}{\textwidth}  % Use \linewidth if nested inside another minipage
\setlength{\tabcolsep}{2pt}  % reduce spacing to avoid overflow
\renewcommand{\arraystretch}{1.1}  % tighten vertical spacing
\rowcolors{6}{}{TableRowAlt} % optional stripe rows

\begin{tabularx}{\textwidth}{
  >{\raggedright\arraybackslash}p{1.2cm}  % Level
  >{\raggedright\arraybackslash}p{1.9cm}  % Proficiency Bonus
  *{11}{>{\raggedright\arraybackslash}X}  % 11 stretch columns
  >{\raggedright\arraybackslash}p{1.2cm}  % Upgrades (final fixed-width)
}
\RedToprule
\multicolumn{13}{l}{\textbf{The Inventor}} \\
\midrule
\multicolumn{6}{c}{\textbf{The Inventor}} & \multicolumn{6}{c}{\textbf{Spell Slots per Spell Level}} & ~ \\
\rowcolor{TableHeader}
\textbf{Level} & \textbf{Proficiency Bonus} & \multicolumn{4}{l}{\textbf{Features}} & \textbf{Spells Known} & \textbf{1st} & \textbf{2nd} & \textbf{3rd} & \textbf{4th} & \textbf{5th} & \textbf{Upgrades} \\
\midrule
1st & +2 & \multicolumn{4}{l}{Inventor Specialization, Magic Item Analysis} & — & — & — & — & — & — & — \\
2nd & +2 & \multicolumn{4}{l}{Spellcasting, Arcane Retrofit, Tool Expertise} & 3 & 2 & — & — & — & — & — \\
3rd & +2 & \multicolumn{4}{l}{Inventor Specialization feature, Specialization Upgrade} & 4 & 3 & — & — & — & — & 1 \\
4th & +2 & \multicolumn{4}{l}{Ability Score Improvement} & 4 & 3 & — & — & — & — & 1 \\
5th & +3 & \multicolumn{4}{l}{Inventor Specialization Feature} & 5 & 4 & 2 & — & — & — & 2 \\
6th & +3 & \multicolumn{4}{l}{Arcane Reconstruction, Cross-Disciplinary Knowledge} & 5 & 4 & 2 & — & — & — & 2 \\
7th & +3 & \multicolumn{4}{l}{Wondrous Item Proficiency} & 6 & 4 & 3 & — & — & — & 3 \\
8th & +3 & \multicolumn{4}{l}{Ability Score Improvement} & 6 & 4 & 3 & — & — & — & 3 \\
9th & +4 & \multicolumn{4}{c}{—} & 7 & 4 & 3 & 2 & — & — & 4 \\
10th & +4 & \multicolumn{4}{l}{Improved Magical Crafting, Wondrous Item Recharge} & 7 & 4 & 3 & 2 & — & — & 4 \\
11th & +4 & \multicolumn{4}{l}{Study of Magic} & 8 & 4 & 3 & 3 & — & — & 5 \\
12th & +4 & \multicolumn{4}{l}{Ability Score Improvement} & 8 & 4 & 3 & 3 & — & — & 5 \\
13th & +5 & \multicolumn{4}{c}{—} & 9 & 4 & 3 & 3 & 1 & — & 6 \\
14th & +5 & \multicolumn{4}{l}{Inventor Specialization Feature} & 9 & 4 & 3 & 3 & 1 & — & 6 \\
15th & +5 & \multicolumn{4}{c}{—} & 10 & 4 & 3 & 3 & 2 & — & 7 \\
16th & +5 & \multicolumn{4}{l}{Ability Score Improvement} & 10 & 4 & 3 & 3 & 2 & — & 7 \\
17th & +6 & \multicolumn{4}{c}{—} & 11 & 4 & 3 & 3 & 2 & 1 & 8 \\
18th & +6 & \multicolumn{4}{l}{Wondrous Item Mastery} & 11 & 4 & 3 & 3 & 3 & 1 & 8 \\
19th & +6 & \multicolumn{4}{l}{Ability Score Improvement} & 12 & 4 & 3 & 3 & 3 & 2 & 9 \\
20th & +6 & \multicolumn{4}{l}{Peerless Inventor} & 12 & 4 & 3 & 3 & 3 & 2 & 9 \\
\bottomrule
\end{tabularx}
\end{minipage}

\vspace{2.5em}
\begin{multicols}{2}
\subsubsection{Specialization Upgrade}

At 3rd level, choose an upgrade from the list at the end of your specialization, gaining the benefits listed in its description. 

You select an additional upgrade at 5th, 7th, 9th, 11th, 13th, 15th, 17th, and 19th level. You can’t select an upgrade more than once, unless the upgrade’s description says otherwise. Whenever you gain a level in this class, you can exchange one of your existing upgrades with another upgrade of the same level requirement as the replaced upgrade.

Whenever an existing upgrade is exchanged for a replacement upgrade (either by a subclass feature or when you gain a level), the new upgrade is selected as if you’re the level you were when you originally gained that upgrade. For example, if you replace your Stormforged Weapon and reselect all of your upgrades as a 5th-level inventor, you could select one 3rd-level upgrade and one 5th-level upgrade, or two 3rd-level upgrades, but you would not be able to select two 5th-level upgrades.
\begin{CalloutBox}{Customizing Inventor Upgrades}
The upgrades for each Inventor Specialization are presented in a list at the end of the subclass, but invariably there will always be ideas for upgrades not included in that list. At the heart of an inventor beats an unrelenting drive for creativity, after all! Feel free to consult your GM for potential custom upgrades.
\end{CalloutBox}
\end{multicols}
\newpage
\twocolumn
The following generic upgrades are available to all subclasses, and can be selected in place of any subclass upgrade:
\begin{fiveitemize}
	\item Shield Proficiency Unrestricted Generic Upgrade 
	\begin{fiveitemize}
		\item You gain proficiency with shields.
	\end{fiveitemize}
	\item Tool Proficiency Unrestricted Generic Upgrade 
	\begin{fiveitemize}
		\item You gain proficiency with a tool of your choice.
	\end{fiveitemize}
\end{fiveitemize}

\subsubsection{Ability Score Improvement}

When you reach 4th level, and again at 8th, 12th, 16th, and 19th level, you can increase one ability score of your choice by 2, or you can increase two ability scores of your choice by 1. As normal, you can’t increase an ability score above 20 using this feature.

\subsubsection{Arcane Reconstruction}

At 6th level, you have mastered the knowledge of using magic to repair things. You learn the mending cantrip, and you can cast it at will. Additionally, you learn the cure wounds spell. If you already know cure wounds, you learn a different spell of your choice from the inventor spell list. Constructs targeted by your cure wounds can regain hit points as normal.

\subsubsection{Cross-Disciplinary Knowledge}

At 6th level, you can expand on your knowledge across fields. You can craft one of the following: any of the Gadgetsmith’s Unrestricted Upgrades, an Infusionsmith’s Animated Weapon, Blasting Rod, or Infused Weapon, a Potionsmith’s Alchemical Reagent Pouch and Alchemical Fire or Alchemical Acid, or a Thundersmith’s Stormforged Weapon.

If you select a Stormforged Weapon, you gain proficiency with that weapon and knowledge of how to make ammunition for it (if required). You can’t apply Infused Weapon to another weapon granted by this class.

If this crafted item is lost or destroyed, you can remake it following the rules for remaking it in its respective subclass.

\subsubsection{Wondrous Item Proficiency}

At 7th level, your familiarity with the workings of magical items means that you ignore all class requirements on the use of magic items. Additionally, you can now attune to up to four magic items at once.

\subsubsection{Improved Magical Crafting}

At 10th level, your experience in creating your own wondrous inventions makes you more adept at crafting a magic item than normal spellcasters. Creating a non-consumable magic item takes you half the time it would normally take.

Additionally, you can make 1 hour of progress toward crafting a magic item, potion, or scroll during a long rest.

\subsubsection{Wondrous Item Recharge}

Starting at 10th level, you can recharge a magic item that has charges or per rest uses, as long as those charges or uses can only be used to cast spells. To restore charges or uses, you perform a ritual that takes 1 minute and expend a spell slot of a level equal to or greater than the spell slot level of a spell cast by the item.

\begin{CalloutBox}{Interaction with the Crafting System}
If you are using the Crafting System presented later in the book, this reduces the time a crafting check takes to 1 hour, allowing you to make two crafting checks during a long rest, or eight during a work day.
\end{CalloutBox}

The number of charges or uses restored to the item is equal to the number of charges or uses required to cast that spell using the item once.

\subsubsection{Study of Magic}

At 11th level, your proficiency in the workings of magic has become so great you can cast the detect magic and identify spells at will without expending a spell slot.

Additionally, you have advantage on Intelligence (Arcana) checks made to understand the workings of magical traps, effects, or runes.

\subsubsection{Wondrous Item Mastery}

Starting at 18th level, you can use a magic item that would normally take an action as a bonus action instead. In addition, you can now attune to up to five magic items at once.

\subsubsection{Peerless Inventor}

At 20th level, your mind is always thinking of new options and clever solutions. At the end of a short or long rest, you can create a temporary version of an Upgrade from your subclass that you meet the prerequisites for but don’t already have. This upgrade must have a level requirement of 11th level or lower, and it lasts until you finish a short or long rest, at which time you can select a temporary Upgrade with this feature again.

\begin{CalloutBox}{Magical or Mechanical?}

One of the core debates that arises around inventors is if they are a “Magical Engineer” or what a magical engineer would even mean. Do they tinker with mechanics and gears or with magic and runes? The answer is, of course, whatever suits your world, game, and the vision of the player. Don't overcomplicate this - we are playing a game with creativity as the backbone foundation!
\end{CalloutBox}
\SectionSplashPage[0.4\textheight]{Gadgetsmith.png}
\begin{multicols}{2}
\section{Gadgetsmith}
A gadgetsmith is an inventor whose curiosity and inventive genius have run rampant. While other inventors may spend their whole career perfecting a single-minded pursuit, a gadgetsmith believes that quantity is at least as good as quality.

Quick footed and even quicker witted, gadgetsmiths are never caught without another trick up their sleeve—their minds always jumping ahead to solve the next problem with a clever contraption.

A gadgetsmith can come from any walk of life, but they usually exemplify curiosity and a distaste for the suppression of knowledge or technology, usually favoring freedom to experiment and leaning toward more chaotic behaviors.
\subsection{Gadgetsmith’s Proficiency}

When you choose this specialization at 1st level, you gain proficiency with nets, rapiers, whips, and tinker’s tools.

\subsection{Essential Tools}

At 1st level, you’ve mastered the creation of the essential reusable tools for surviving the battlefield as a gadgeteer. You have the following items:

\begin{fiveitemize}
	\item \textbf{Gadgetsmith Weapon.} \\ Choose one of the following from the Upgrade section: 
	\begin{fiveitemize}
		\item Impact Gauntlets, Lightning Baton, Repeating Hand Crossbow, Ricocheting Weapon, or Shock Generator. 
		\item You gain the chosen upgrade, and it doesn’t count against your total number of upgrades.
	\end{fiveitemize}
	\item \textbf{Grappling Hook.} \\ You can forgo any attack you make as part of the Attack action to throw your grappling hook at a target within 20 feet of you. The target can be a creature, object, or surface. If the target is Small or smaller, make a grapple check against it, pulling it to you and grappling it on a success. If the target is Medium or larger, you can instead pull yourself to it (no check required); this doesn’t grapple the target. Opportunity attacks provoked by this movement are made with disadvantage.
\newpage
\end{fiveitemize}
\end{multicols}
\twocolumn
\begin{fiveitemize}
	\item \textbf{Smoke Bomb.} \\ As an action, you can use a smoke bomb to cast the fog cloud spell centered on yourself, without expending a spell slot and requiring no components. The spell lasts for a number of rounds equal to your Intelligence modifier and doesn’t require your concentration. When you cast fog cloud in this way, you can choose for it to have a radius of 5, 10, 15, or 20 feet.
\end{fiveitemize}

\begin{CalloutBox}{Unlimited Smoke Bombs}
If you or your GM are concerned about the unlimited use of Smoke Bombs, consider limiting them to a number of uses per short rest equal to your proficiency bonus. This shouldn’t mechanically impact their usage, but the reason they don’t have a limit is to encourage players to feel free to use them in non-combat situations. With limits—even ones that are not usually reached—players tend to treat their Smoke Bomb as an exclusively combat-oriented ability. Alternatively, if you are using the Crafting Mechanics introduced later in this book, you can require the player to keep up with them as a type of ammunition crafted.
\end{CalloutBox}
\subsection{Additional Upgrade}

At 3rd level, you’ve mastered the essential tools and have begun to tinker with ways to expand your arsenal. The number of upgrades you can have, based on your class level, increases by one. It increases by one again at 5th level (for a total of two more than shown in the Inventor class table).

\subsection{Recycle Gadgets}

Starting at 3rd level, you can disassemble your gadgets and create different ones during a long rest. When you finish the rest, you can replace any upgrade you have with a new one.

The chosen upgrade must still be one that is valid for the level at which you gained the replaced upgrade. For example, at 9th level, you can only have one upgrade that has a 9th-level prerequisite.

Additionally, if one of your gadgets is destroyed, you can use this feature to recreate it over the course of a long rest, which requires materials worth 20 gp.

\subsection{Extra Attack}

Beginning at 5th level, you can attack twice, instead of once, whenever you take the Attack action on your turn.

\subsection{Combat Gadgets}

Beginning at 14th level, you can forgo any attack you make as part of the Attack action to use a gadget that normally requires an action to use.

\section{Gadgetsmith Upgrades}
\subsection{Unrestricted Upgrades}
\begin{fiveitemize}

	\item \textbf{Airburst Mine.}
	
	You create a mechanical device capable of producing a devastating blast. You can use this device to cast shatter or thunderburst mine* without expending a spell slot.

	Once used, this gadget can’t be used again until you finish a short or long rest.

	\item \textbf{Belt of Adjusting Size.}
	
	You create a belt with a creature size dial on it. While you are wearing the belt, you can use an action to cast the enlarge/ reduce spell on yourself without expending a spell slot.

	Once used, this gadget can’t be used again until you finish a short or long rest.

	\item \textbf{Element Eater.}
	
	You create a device capable of absorbing incoming elemental damage. You can activate this device as a reaction when you take acid, cold, fire, lightning or thunder damage. The device grants you resistance to the triggering damage type until the start of your next turn. After triggering it, the first time you hit with a melee attack on your next turn, the target takes an extra 1d6 elemental damage of the triggering type.

	Once used, this gadget can’t be used again until you finish a short or long rest.

	\item \textbf{Enhanced Grappling Hook.}
	
	You enhance the grappling hook gained from the Essential Tools feature, increasing its range to 30 feet. Additionally, when pulling yourself to a Large or larger target, the enhanced power of the grappling hook allows you to drag one Medium or smaller creature with you. The creature must be willing or grappled by you, and it must be within 5 feet of you.

	\item \textbf{Fire Spitter.}
	
	You create a gadget that can create a quick blast of fire. As an action, you can use the gadget to cast burning hands as a 2nd-level spell without expending a spell slot.

	Once used, this gadget can’t be used again until you finish a short or long rest.

	\item \textbf{Flashbang.}
	
	You create a high-luminosity discharge device. As an action, you can use the device and target a point within 30 feet of you. Any creature within 20 feet of that point must succeed on a Dexterity saving throw against your spell save DC or be blinded until the end of its next turn.

	Once used, this gadget can’t be used again until you finish a short or long rest.

	\item \textbf{Gliding Cloak.}
	
	You make a cloak that allows you to glide while falling. When you fall more than 10 feet and aren’t incapacitated, you can spread the cloak to reduce your rate of descent to 30 feet per round, taking no falling damage when you land. While falling in this manner under normal gravity, you can move up to 2 feet horizontally for every 1 foot you descend.

	\item \textbf{Gravity Switch.}
	
	You build a switch that turns off gravity. You can use it to cast fall* without expending a spell slot.

	Once used, this gadget can’t be used again until you finish a short or long rest.

	\item \textbf{Impact Gauntlet.}
	
	You create a magic weapon capable of amplifying the impact of your blows. You have proficiency with this weapon, and it has the finesse and light properties, as well as the special property (described below). It deals 1d6 bludgeoning damage on a hit. You can select this upgrade up to two times, making a separate weapon each time.

	\textit{Special.} When you make an attack with this weapon, you can choose to forgo adding your proficiency bonus to the attack roll. If your attack hits, you can add double your proficiency bonus to the damage roll.
	\begin{fiveitemize}
	
		\item \textbf{Variant: Power Fist.}

		Your GM may allow you to take the Power Fist upgrade from the Warsmith specialization in place of an Impact Gauntlet, which is similar but deals 1d8 damage and lacks the finesse property. For the purpose of upgrades, Impact Gauntlets and Power Fists should be considered interchangeable.
	\end{fiveitemize}
	
	\item \textbf{Jumper Cable (Prerequisite: Shock Generator).}
	
	Once per turn when you deal lightning damage with a cantrip from your Shock Generator, you can add your Intelligence modifier to the damage dealt.

	Additionally, you can use your Shock Generator as an action to make a DC 10 Intelligence (Medicine) check, attempting to stabilize a dying creature within your reach. If you succeed on the check, the creature regains 1 hit point, becoming conscious again, and gains a number of temporary hit points equal to your inventor level. When you stabilize a creature in this way, it gains one level of exhaustion.

	\item \textbf{Jumping Boots.}
	
	You modify your boots with arcane boosters. While wearing these boots, you are under the effects of the jump spell.

	\item \textbf{Lightning Baton.}
	
	You create a magic weapon that channels lightning. You have proficiency with this weapon, and it has the finesse and light properties. It deals 1d4 bludgeoning damage plus 1d4 lightning damage on a hit. When you roll a 20 on an attack roll made with this weapon, the target must succeed a Constitution saving throw against your spell save DC or become stunned until the start of your next turn. You can choose this upgrade up to two times, making a separate weapon each time.

	\item \textbf{Mechanical Arm.}
	
	You create a mechanical arm, giving you an extra hand. This mechanical arm only functions while it is mounted on gear you are wearing, but can be operated mentally without the need for your hands. This mechanical arm can serve any function a normal hand could, such as holding things, making attacks, and interacting with the environment, but it doesn’t grant you any additional actions.

	\item \textbf{Mechanical Familiar.}
	
	You create the blueprint for a small, mechanical creature. At the end of a long rest, you can choose to create a familiar based on it, casting the find familiar spell without expending a spell slot. When cast in this way, the familiar is a construct. It stays active until you deactivate it or it is destroyed.

	In either case, you can choose to reactivate it at the end of a long rest.

	\item \textbf{Net Launcher.}
	
	You build a device capable of delivering nets to their targets more effectively. While you have this upgrade, nets have a normal range of 20 feet and a long range of 60 feet for you.

	\item \textbf{Quick Essential Gadget.}
	
	You modify your essential gear for quickened use. You can use your Grappling Hook or Smoke Bomb as a bonus action.

	Once you use either of them in this way, you can’t use either of them in this way again until you finish a short or long rest.

	\item \textbf{Repeating Hand Crossbow.}
	
	You build an improved hand crossbow. You have proficiency with this weapon, which has a normal range of 30 feet and a long range of 120 feet, and it has the ammunition and light properties, as well as the special property (described below). It deals 1d6 piercing damage on a hit.
	
	\textit{Special.} This weapon doesn’t require a free hand to load, as it has a built-in loading mechanism. When you make an attack with this weapon as part of the Attack action, and you don’t have disadvantage on the attack roll, you can choose to make the attack roll with disadvantage to make one additional attack with this weapon as a bonus action. The attack roll for this additional attack is also made with disadvantage.
	
	\item \textbf{Ricocheting Weapon.}
	
	You create a thrown melee weapon enchanted or engineered to ricochet off its targets and return to your hand. You have proficiency with this weapon, which has a normal range of 30 feet and a long range of 90 feet, and it has the finesse and thrown properties, as well as the special property (described below). It deals 1d8 bludgeoning, piercing, or slashing damage on a hit (your choice when you choose this upgrade). You can Arcane Retrofit this weapon.


	\textit{Special.} When this weapon is thrown, you can target two creatures within 10 feet of each other, making a separate attack roll against each target; the damage dealt by the attack is halved for each target hit after the first. In addition, the weapon flies back to your hand immediately after the attack.

	\item \textbf{Shock Generator.}
	
	You create a device capable of generating potent shocks. You can use this device to cast shocking grasp . When you cast shocking grasp in this way, you can use either your Dexterity or Intelligence modifier for the attack roll.

	\item \textbf{Shocking Hook (Prerequisite: Shock Generator).}
	
	You can integrate your Shock Generator and your Grappling Hook. Immediately after using your Grappling Hook to pull a creature to you or to pull yourself to a creature, you can use your bonus action to cast shocking grasp from your Shock Generator targeting that creature.

	\item \textbf{Sight Lenses.}
	
	You create a set of lenses that allow you to see through darkness and obscurement, which you can integrate into a pair of goggles, glasses, or other vision assistance. While using these lenses, you can see through fog, mist, smoke, clouds, and nonmagical darkness out to a range of 15 feet.

	\item \textbf{Smoke Cloak.}
	
	You create a cloak that causes you to blend in with smoke. When you are lightly or heavily obscured by smoke at the start of your turn, you are invisible until your turn ends, or until you cast a spell, make an attack, or deal damage to a creature.

	\item \textbf{Steelweave Nets.}
	
	You thread your nets with metal reinforcement, making them tougher and more conductive. The net has an AC of 15 and resistance to slashing damage. Whenever a creature restrained by the net takes lightning damage, it takes an extra 1d6 lightning damage.
	Only you have proficiency with your Steelweave Nets. You can have a number of these nets up to your Intelligence modifier at a time (a minimum of one), and you can replenish your supply of them over the course of a long rest by forging normal nets into Steelweave Nets.

	\item \textbf{Striding Boots.}
	
	You modify your boots with amplified striding speed. While wearing these boots, you are under the effects of the longstrider spell.

	\item \textbf{Trick Shots.}
	
	At the end of a short or long rest, you can add one of the following effects to a piece of ammunition or a thrown weapon, turning it into a trick shot; the effect is expended when you use the trick shot:
	\begin{fiveitemize}
		\item \textit{Arcane Trick.} You can imbue a cantrip or 1st-level spell you know into the shot, casting it on impact where the shot lands. The target of the attack is also the spell’s target.
		\item \textit{Bouncing.} You can attack a target within range that has total cover from you if you know the target’s location, bouncing the shot off a surface. The target has the benefit of half cover for the attack, instead of total cover.
		\item \textit{Ricochet.} If you hit a target with the shot, you can make another attack roll against a different target within 10 feet of the first as part of the same attack.
		\item \textit{Smoke Shot.} You can use the shot to deploy a Smoke Bomb on impact where the shot lands.
		\item \textit{Special Tip.} The shot deals one of the following damage types of your choice: bludgeoning, piercing, slashing, acid, cold, fire, or lightning.
	\end{fiveitemize}
	You can select this upgrade multiple times, allowing you to make multiple trick shots per rest. You can have a maximum number of trick shots equal to the number of times you have selected this upgrade.
\end{fiveitemize}
\subsection{5th-Level Upgrades}
\begin{fiveitemize}

	\item \textbf{Antimagical Shackle.}
	
	You create a set of antimagical shackles. As an action when you are adjacent to a creature, you can attempt to shackle the creature to yourself or a nearby object using these shackles. Make a Dexterity (Sleight of Hand) check contested by the target’s Strength (Athletics) or Dexterity (Acrobatics) check (target’s choice). A creature that is immune to being grappled or restrained automatically succeeds on this check. If you succeed, the target is shackled to you or the nearby object, and it can only move if it’s able to move what it’s shackled to.

	Additionally, while shackled in this way, the target can’t teleport, travel to another plane of existence, alter its form (such as by the polymorph spell), or dematerialize. As an action, the target can make a Strength saving throw against your spell save DC, breaking the shackles on a success; otherwise, the shackles last until you remove them.

	\item \textbf{Binding Rope.}
	
	You create a rope that is capable of animating and binding a creature. As an action, choose a creature within 30 feet of you. The target must make a Dexterity saving throw against your spell save DC or become restrained until the end of your next turn. If you are currently grappling the target, it makes the Dexterity saving throw with disadvantage. The rope can only restrain one target at a time.

	\item \textbf{Crossbow Spider.}
	
	You modify a crossbow to be able to aim and fire remotely. As an action, you can deploy a Tiny construct called a crossbow spider (see the construct’s statistics below). Once deployed, you can use a bonus action to cause the construct to move up to its speed and make a crossbow attack from its location. The construct becomes inactive after 1 minute, or after it has been used to make 10 attacks.

	Once you have activated the construct, you can’t activate it again until you finish a long rest, unless you expend a spell slot of 1st level or higher to activate it again. If the construct is destroyed, you can’t activate it again until you repair it or recreate it during a long rest.

	The construct can’t take any actions besides the ones you direct it to take with your bonus action.
\end{fiveitemize}
\begin{StatblockBackground}

\noindent\raggedright{\color{DndRed}\Large\textbf{Crossbow Spider}}\\
\raggedright\textit{Tiny Construct, unaligned}
\vspace{2pt}
\MonsterGradientRule
\vspace{2pt}

\begin{tabularx}{\columnwidth}{>{\raggedright\arraybackslash}p{0.33\columnwidth} X}
\textbf{\color{DndRed}Armor Class} & 10 \\
\textbf{\color{DndRed}Hit Points} & 5 (1d6 + 2) \\
\textbf{\color{DndRed}Speed} & 15 ft., climb 15 ft. \\
\end{tabularx}

\vspace{2pt}
\MonsterGradientRule
\vspace{2pt}

\begin{tabularx}{\columnwidth}{*{6}{>{\centering\arraybackslash}X}}
\textbf{\color{DndRed}STR} & \textbf{\color{DndRed}DEX} & \textbf{\color{DndRed}CON} &
\textbf{\color{DndRed}INT} & \textbf{\color{DndRed}WIS} & \textbf{\color{DndRed}CHA} \\
2 (−4) & 10 (+0) & 10 (+0) & 1 (−5) & 1 (−5) & 1 (−5) \\
\end{tabularx}

\vspace{2pt}
\MonsterGradientRule
\vspace{2pt}

\begin{tabularx}{\columnwidth}{>{\raggedright\arraybackslash}p{0.33\columnwidth} X}
\textbf{\color{DndRed}Damage Immunities} & poison \\
\textbf{\color{DndRed}Condition Immunities} & charmed, exhaustion, frightened, paralyzed, petrified, poisoned \\
\textbf{\color{DndRed}Senses} & passive Perception 5 \\
\textbf{\color{DndRed}Languages} & — \\
\textbf{\color{DndRed}Proficiency Bonus} & equals your bonus \\
\textbf{\color{DndRed}Challenge} & — \\
\end{tabularx}

\vspace{2pt}
\MonsterGradientRule
\vspace{2pt}

\MonsterSection{Actions}


\MonsterAction{Crossbow}{Ranged Weapon Attack: your spell attack modifier to hit, range 30/120 ft., one target you can see. Hit: 1d6 + your Intelligence modifier piercing damage.}

\end{StatblockBackground}
\begin{fiveitemize}

	\item \textbf{Explosive Gauntlet.}
	
	Your gauntlets can exert massive kinetic force when striking. When you make an attack with your Impact Gauntlet, you can choose to be pushed 10 feet in the opposite direction of your attack. Alternatively, you can use a bonus action to attempt to push the target after the attack. The target must succeed on a Strength saving throw or be pushed 10 feet away from you. This movement doesn’t provoke opportunity attacks.

	You can use the movement effect of this gauntlet even if you’re not attacking a target, allowing you to push yourself in any direction (including upwards).

	\item \textbf{Smoky Images.}
	
	Immediately after using a Smoke Bomb, you can use a bonus action to cast mirror image without expending a spell slot.

	Once you cast the spell in this way, you can’t do so again until you finish a short or long rest.

	\item \textbf{Vanishing Trick.}
	
	Immediately after using a Smoke Bomb, you can cast misty step without expending a spell slot.

	Once you cast the spell in this way, you can’t do so again until you finish a short or long rest.
\end{fiveitemize}

\subsection{9th-Level Upgrades}
\begin{fiveitemize}

	\item \textbf{Arcane Nullifier.}
	
	You make a device that nullifies the arcane through means you assure everyone else you understand. As an action, you can use this device to cast dispel magic without expending a spell slot.

	Once used, this gadget can’t be used again until you finish a short or long rest.

	\item \textbf{Phase Trinket.}
	
	You create a magical gadget that manipulates ethereal magic. As an action, you can use this device to cast blink or dimension door without expending a spell slot.

	Once used, this gadget can’t be used again until you finish a long rest.

	\item \textbf{Stinking Gas.}
	
	You make a more potent compound for your Smoke Bombs. When you use a Smoke Bomb, you can choose to cast stinking cloud , instead of fog cloud , following the same rules.

	Once you cast stinking cloud in this way, you can’t do so again until you finish a short or long rest.

	\item \textbf{Stopwatch Trinket.}
	
	You create a magical stopwatch that manipulates time magic. As an action, you can use the stopwatch to cast haste or slow without expending a spell slot.

	Once used, this gadget can’t be used again until you finish a long rest.
\end{fiveitemize}

\subsection{11th-Level Upgrades}
\begin{fiveitemize}

	\item \textbf{Flying Gadget.}
	
	You build a device that allows you to fly, such as deployable artificial wings. You can activate this device as a bonus action, or as a reaction when you fall. Whatever form the device takes, when it’s activated, you gain a flying speed of 30 feet, which lasts until you choose to deactivate it (no action required).

	\item \textbf{Lightning Generator (Prerequisite: Shock Generator).}
	
	You upgrade your Shock Generator with additional lightning capabilities. You can overload your Shock Generator to cast lightning bolt without expending a spell slot.

	Once you cast lightning bolt in this way, you can’t do so again until you finish a short or long rest.

	Additionally, once per turn when you deal lightning damage with a spell, you can add your Intelligence modifier to the spell’s damage roll.

	\item \textbf{Truesight Lenses (Prerequisite: Sight Lenses).}
	
	You upgrade and fine-tune your Sight Lenses, granting you truesight out to a range of 15 feet while using them.

	\item \textbf{Useful Universal Key.}
	
	You create a Universal Key to obstacles, transmuting them into not-obstacles. As an action, you can apply this key to a surface to cast knock or passwall on it without expending a spell slot.

	Once used, this gadget can’t be used again until you finish a long rest.
\end{fiveitemize}

\subsection{15th-Level Upgrades}
\begin{fiveitemize}

	\item \textbf{Bee Swarm Rockets.}
	
	You design a type of tiny, firecracker-like device, which can release rockets in large numbers. You have a maximum number of rockets equal to your inventor level. As an action, you can release between 1 and 10 rockets. Each rocket targets a point you can see within 40 feet of you. Each creature within 10 feet of one of these points must make a Dexterity saving throw against your spell save DC.

	Creatures in the area of effect of multiple rockets make a separate saving throw for each rocket. For each failed save, a creature takes 2d6 + 1 fire damage, or half as much damage on a successful one.

	You can replenish your supply of rockets to your maximum over the course of a long rest.

	\item \textbf{Bracers of Empowerment.}
	
	You create bracers that can empower you. You can use these bracers to cast martial transformation* without expending a spell slot.

	Once used, this gadget can’t be used again until you finish a long rest.

	\item \textbf{Dimensional Toolbox.}
	
	You build a toolbox, filling it with the many ideas you’ve had and discarded, with the magical power of making those ideas reality when you need them most. As an action, you can reach into the toolbox and withdraw an Unrestricted Upgrade of your choice (one with no level requirement) from the Gadgetsmith list. Gadgets withdrawn from the toolbox are fleeting, and disappear after 1 minute.

	Once you have withdrawn an upgrade from the toolbox, you can’t withdraw another upgrade from it until you finish a long rest.

	\item \textbf{Disintegration Ray.}
	
	You create a Disintegration Ray. You can use this device to cast disintegrate without expending a spell slot.

	Once used, this gadget can’t be used again until you finish a long rest.

	\item \textbf{Nexus Hive.}
	
	You create a floating metal sphere, that when activated spews a storm of metal shards in all directions around it. As an action you can deploy a small construct (see the construct’s statistics below), throwing it up to 30 feet where it will hover in place. Once deployed you can use a bonus action to move the Nexus Hive up to 20 ft. The construct becomes inactive after 1 minute.
	If the construct is destroyed, you can’t activate it again until you repair it or recreate it during a long rest.

	The construct can’t take any actions besides the ones you direct it to take with your bonus action.
\end{fiveitemize}

\begin{StatblockBackground}

\noindent\raggedright{\color{DndRed}\Large\textbf{Nexus Hive}}\\
\raggedright\textit{Small Construct, unaligned}
\vspace{2pt}
\MonsterGradientRule
\vspace{2pt}

\begin{tabularx}{\columnwidth}{>{\raggedright\arraybackslash}p{0.33\columnwidth} X}
\textbf{\color{DndRed}Armor Class} & 15 (natural armor) \\
\textbf{\color{DndRed}Hit Points} & 30 \\
\textbf{\color{DndRed}Speed} & 0 ft., fly 20 ft. (hover) \\
\end{tabularx}

\vspace{2pt}
\MonsterGradientRule
\vspace{2pt}

\begin{tabularx}{\columnwidth}{*{6}{>{\centering\arraybackslash}X}}
\textbf{\color{DndRed}STR} & \textbf{\color{DndRed}DEX} & \textbf{\color{DndRed}CON} &
\textbf{\color{DndRed}INT} & \textbf{\color{DndRed}WIS} & \textbf{\color{DndRed}CHA} \\
2 (−4) & 10 (+0) & 12 (+1) & 1 (−5) & 1 (−5) & 1 (−5) \\
\end{tabularx}

\vspace{2pt}
\MonsterGradientRule
\vspace{2pt}

\begin{tabularx}{\columnwidth}{>{\raggedright\arraybackslash}p{0.33\columnwidth} X}
\textbf{\color{DndRed}Damage Immunities} & poison, psychic \\
\textbf{\color{DndRed}Condition Immunities} & charmed, exhaustion, frightened, paralyzed, petrified, poisoned \\
\textbf{\color{DndRed}Senses} & passive Perception 5 \\
\textbf{\color{DndRed}Languages} & — \\
\textbf{\color{DndRed}Proficiency Bonus} & equals your bonus \\
\textbf{\color{DndRed}Challenge} & — \\
\end{tabularx}

\vspace{2pt}
\MonsterGradientRule
\vspace{2pt}

\MonsterSection{Traits}

\MonsterTrait{Swarm.}{The nexus hive sprays metal shards in a 10-foot radius around itself. A creature takes 4d4 piercing damage when it enters this area for the first time on a turn or starts its turn there.}
\MonsterTrait{Explosive Vent.}{When reduced to 0 hit points, the nexus hive will violently vent energy. Creatures within 10 feet of the nexus hive must succeed on a Dexterity saving throw against your spell save DC. A creature takes 4d6 fire damage on a failed save, or half as much damage on a successful one.}


\end{StatblockBackground}

\SectionSplashPage[0.4\textheight]{Golemsmith.png}
\begin{multicols}{2}
\section{Golemsmith}
Golemsmiths are inventors who have committed themselves to creating a true work of artifice, forging a golem. With this painstaking life ambition, they plan and design meticulously, even if, in practice, sometimes compromises on materials must be made.

Why a golemsmith embarks on the quest to forge this artificial construct of life can vary. For many it is the pure pursuit of forging the perfect creation; while for others, it is simply so they do not have to carry around their loot, or to have a loyal companion to count on at all times.

A golemsmith is rarely chaotic, as they are people of great care and discipline—those that are not would not have succeeded where they have. However, some have been set on their path by such events that might drive them to interact chaotically with society as a whole.
\subsection{Golemsmith's Proficiency}
When you choose this specialization at 1st level, you gain proficiency with smith’s tools and tinker’s tools.
\subsection{Mechanical Golem}
Starting at 1st level, you forge a mechanical golem to carry out your orders and protect you. The golem is under your control and understands the languages you speak, but it can’t speak.

The golem obeys your commands as best it can. On your turn, you can verbally command it where to move (no action required by you), and you can use your action to command it to take an action. Additionally, whenever it would be able to take a reaction, you must use your reaction to command it to take that reaction. The golem acts on your commands during your turn. If you issue no commands to your golem, it takes no actions.

\textbf{The base statistics for your golem are as follows:}

\begin{fiveitemize}
	\item Your golem’s proficiency bonus increases when yours does. If the golem is killed, it can be returned to life by normal means, such as the revivify spell, or you can repair it during a long rest if you have access to its body. At the end of the long rest, it returns to life with half its hit points (rounded down).

	\item If the golem is beyond recovery, you can recreate it exactly as it was, with four days of work (8 hours each day) and 100 gp of raw materials.
\end{fiveitemize}
\end{multicols}
\newpage
\twocolumn
\begin{fiveitemize}
	\item Over the course of a short rest, you can restore hit points to your golem equal to your Intelligence modifier + your inventor level, or you can fully repair it to during a long rest, causing it to regain all its hit points. Any spells you cast that cause the target to regain hit points ignore any restrictions against constructs for your golem.
\end{fiveitemize}
\begin{StatblockBackground}

\noindent\raggedright{\color{DndRed}\Large\textbf{Mechanical Golem}}\\
\raggedright\textit{Medium Construct, unaligned}
\vspace{2pt}
\MonsterGradientRule
\vspace{2pt}

\begin{tabularx}{\columnwidth}{>{\raggedright\arraybackslash}p{0.33\columnwidth} X}
\textbf{\color{DndRed}Armor Class} & 14 (natural armor) \\
\textbf{\color{DndRed}Hit Points} & (5 + (the golem's Constitution modifier +5) * Inventor level \\
\textbf{\color{DndRed}Speed} & 30 ft. \\
\end{tabularx}

\vspace{2pt}
\MonsterGradientRule
\vspace{2pt}

\begin{tabularx}{\columnwidth}{*{6}{>{\centering\arraybackslash}X}}
\textbf{\color{DndRed}STR} & \textbf{\color{DndRed}DEX} & \textbf{\color{DndRed}CON} &
\textbf{\color{DndRed}INT} & \textbf{\color{DndRed}WIS} & \textbf{\color{DndRed}CHA} \\
14 (+2) & 12 (+1) & 12 (+1) & 4 (−3) & 5 (−3) & 1 (−5) \\
\end{tabularx}

\vspace{2pt}
\MonsterGradientRule
\vspace{2pt}

\begin{tabularx}{\columnwidth}{>{\raggedright\arraybackslash}p{0.33\columnwidth} X}
\textbf{\color{DndRed}Damage Immunities} & poison, psychic \\
\textbf{\color{DndRed}Condition Immunities} & charmed, exhaustion, frightened, poisoned \\
\textbf{\color{DndRed}Senses} & passive Perception 7 \\
\textbf{\color{DndRed}Languages} & understands the languages of its creator but can't speak. \\
\textbf{\color{DndRed}Proficiency Bonus} & equals your bonus \\
\textbf{\color{DndRed}Challenge} & — \\
\end{tabularx}

\vspace{2pt}
\MonsterGradientRule
\vspace{2pt}

\MonsterSection{Traits}

\MonsterTrait{Bound.}{The golem is magically bound to its creator. As long as the creator and it are on the same plane of existence, the creator can telepathically call the golem to travel to it, and the golem knows the distance and direction to its creator.}

\end{StatblockBackground}

\subsubsection{Golem Chassis}
When you create your golem, you can choose one of the chassis below, adding the prefix to its type. For example, if you select Ironwrought, your golem becomes an Ironwrought Golem and gains the corresponding modifications to its base statistics.

\Subheading{Launcher}

Your golem becomes akin to a mobile turret, taking the frame of a ballista or other launching device. Its Dexterity becomes 16 (+3), its speed becomes 25 feet, and it has the following natural weapon attack, which uses its Dexterity for the attack and damage rolls:
\begin{fiveitemize}
	\item \textbf{Shoot.} \textit{Ranged Weapon Attack:} +5 to hit, range 60/240 ft., one target. Hit: 1d10 + 3 piercing damage.
\end{fiveitemize}

\Subheading{Quadrupedal}

Your golem takes on a quadrupedal design. Larger and sturdier, it is more suitable to launch into the fray or carry its creator. Its size becomes Large, its base AC becomes 16, and its speed becomes 35 feet. In addition, its Strength and Constitution become 16 (+3), and it has the following natural weapon attack, which uses its Strength for the attack and damage rolls:
\begin{fiveitemize}
	\item \textbf{Bite.} \textit{Melee Weapon Attack:} +5 to hit, reach 5 ft., one target. Hit: 1d10 + 3 piercing damage.
\end{fiveitemize}

\Subheading{Specialized}

Your golem defies all expectations, its design fueled by your own rampant creativity. A little less robust and stable, it is far more extensible to your visionary plans. Your golem starts with the basic golem statistics, but you can select two free upgrades for it that don’t count against your upgrade total. In addition, it has the following natural weapon attack, which uses its Strength for the attack and damage rolls:
\begin{fiveitemize}
	\item \textbf{Slam.} \textit{Melee Weapon Attack:} +4 to hit, reach 5 ft., one target. Hit: 1d8 + 2 bludgeoning damage.
\end{fiveitemize}

\Subheading{Optional: Flesh}
When a golemsmith is selecting their golem, some disturbed minds think “but what if I want it be a fleshy monstrosity against sanity?” Well, as usual, the inventor is here for you:

Your golem is roughly humanoid, a lumbering terror of stitched flesh with an odd smell. Its Strength becomes 16 (+3), it gains proficiency with shields, simple weapons, and martial weapons, and it has the following natural weapon attack, which uses its Strength for the attack and damage rolls:
\begin{fiveitemize}
	\item \textbf{Slam.} \textit{Melee Weapon Attack:} +5 to hit, reach 5 ft., one target. Hit: 1d4 + 3 bludgeoning damage.
\end{fiveitemize}

\begin{CalloutBox}{Additional Changes}

Due to your unnatural development as a Fleshsmith, your base Golemsmith subclass features change. Replace your smith’s tools proficiency with leatherworker’s tools proficiency. The upgrades are the same, but with “...but from some gross, fleshy organ.” At your GM’s discretion, you can apply Fleshsmith upgrades to your golem. This has the potential to cause balance issues, so assume the GM will say no unless it’s not too powerful and you can explain how it’s not.
\end{CalloutBox}

\Subheading{Ironwrought}

Your golem is roughly humanoid and comes with the robust flexibility and options that this form provides. Its Strength becomes 16 (+3), it gains proficiency with shields, simple weapons, and martial weapons, and it has the following natural weapon attack, which uses its Strength for the attack and damage rolls:
\begin{fiveitemize}
	\item \textbf{Slam.} \textit{Melee Weapon Attack:} +5 to hit, reach 5 ft., one target. Hit: 1d4 + 3 bludgeoning damage.
\end{fiveitemize}

\Subheading{Winged}

Your golem is modeled off a flying creature. Smaller and more lightweight than most golems, it is kept aloft by intricate wings. Its size becomes Small, it gains a flying speed of 30 feet, and it has the following natural weapon attack, which uses its Strength for the attack and damage rolls:
\begin{fiveitemize}
	\item \textbf{Talon.} \textit{Melee Weapon Attack:} +4 to hit, reach 5ft., one target. Hit: 1d8 + 2 slashing damage.
\end{fiveitemize}

\subsection{Intelligent Oversight}

Starting at 3rd level, you can take the Help action as a bonus action when assisting your golem.

Additionally, when you take the Help action to aid an ally (including your golem) in attacking a creature, the target of that attack can be up to 30 feet away from you, rather than within 5 feet of you, if your ally can see or hear you.

\subsection{Autonomous Action}

Starting at 5th level, you no longer need to spend your action or reaction to command your golem to use its action or reaction, and you can issue commands to it mentally while it is within 60 feet of you.

If the golem isn’t commanded to take any action, it will take the Dodge action in combat.

\subsection{Magical Nature}

Additionally at 5th level, your golem’s natural weapons count as magical for the purpose of overcoming resistance and immunity to nonmagical attacks and damage.

\subsection{Perfected Design}

Starting at 14th level, your golem can add your Intelligence modifier to all of its attack rolls, ability checks, and saving throws.

\section{Golemsmith Upgrades}
\subsection{Unrestricted Upgrades}

\Subheading{Arcane Barrage Armament}
 
You install a mounted armament on your golem, taking whatever form is most appropriate, which is charged with arcane power. As an action, the golem can cast magic missile through the armament, without expending a spell slot and requiring no components.
Once the golem casts the spell in this way, it can’t do so again until you finish a short or long rest.

When you reach 5th level in this class, the spell is cast as a 2nd-level spell, and it increases again when you reach 11th level (as a 3rd-level spell) and 17th level (as a 4th-level spell).

\Subheading{Arcane Resonance}
 
You craft a magical essence connector and install it into your golem’s core, allowing you and the golem to share certain magical effects. You can make any spell you cast that targets only you also target your golem.

\Subheading{Defender Protocol}
 
You build a protocol into your golem, instructing it to defend its master. The golem gains the Protection Fighting Style.

\Subheading{Ether Heart}

You install a magical ether heart into your golem, because that seems like a good idea. Your golem gains 2 charges. It can expend 1 charge to cast any 1st-level spell you know, and it regains all expended charges when you finish a long rest.

\Subheading{Flamethrower Armament}
 
You install an armament to your golem, taking whatever form is most appropriate, which is capable of producing powerful flames. As an action, the golem can cast burning hands through the armament, without expending a spell slot and requiring no components. The spell save DC is equal to your spell save DC.

Once the golem casts the spell in this way, it can’t do so again until you finish a short or long rest.

When you reach 5th level in this class, the spell is cast as a 2nd-level spell, and it increases again when you reach 11th level (as a 3rd-level spell) and 17th level (as a 4th-level spell).

\Subheading{Fine-Tuned Dexterity}
 
You craft improved gears and joints for your golem. Your golem’s Dexterity score increases by 2, increasing its ability to perform tasks that require fine motor skills. If this upgrade increases the golem’s Dexterity to 16 or higher, it gains proficiency with thieves’ tools; if this upgrade increases its Dexterity to 18, it gains proficiency in the Stealth skill.

You can select this upgrade multiple times, but you can’t increase the golem’s Dexterity above 18 using this upgrade.

\Subheading{Grappling Appendages}
 
You install an additional pair of grappling appendages to your golem, which take a form of your choice. For the purpose of grappling, these appendages count as two additional free hands for your golem. Additionally, as long as your golem isn’t using these appendages to grapple, it has a climbing speed equal to its walking speed.

\Subheading{Heavy Armor Plating}
 
You can incorporate a suit of heavy armor into your golem, allowing it to calculate its armor class as if it was wearing that armor. Your golem is proficient with the incorporated armor, and it has disadvantage on Dexterity (Stealth) checks until the armor is removed. You can remove the incorporated armor or incorporate a new suit of armor at any time, but it takes twice as long as it would normally take to don or doff the armor.

\Subheading{Magical Essence}
 
You infuse a fragment of magical essence into your golem, allowing it to attune to one magic item. This follows all normal attunement rules.

\Subheading{Structural Constitution}
 
You have reinforced your golem with layers of protection and redundant systems. Your golem’s Constitution score increases by 2, increasing its stability and durability. If this upgrade increases the golem’s Constitution to 16 or higher, it gains proficiency in Constitution saving throws; if this upgrade increases its Constitution to 18, it has advantage on death saving throws.

You can select this upgrade multiple times, but you can’t increase the golem’s Constitution above 18 using this upgrade.

\Subheading{Systematic Strength}
 
You build an improved frame and power source for your golem. Your golem’s Strength score increases by 2, increasing its ability to perform tasks that require raw strength. If this upgrade increases the golem’s Strength to 16 or higher, it gains proficiency in Strength saving throws; if this upgrade increases its Strength to 18, it gains proficiency in the Athletics skill.

You can select this upgrade multiple times, but you can’t increase the golem’s Strength above 18 using this upgrade.

\Subheading{Warfare Routines}
 
You advance the control routines for your golem, allowing it to fight more effectively. Your golem gains one of the following Fighting Styles of your choice: Archery, Dueling, or Great Weapon Fighting.

\subsection{5th-Level Upgrades}

\Subheading{Cloaking Device} 

You install an arcane cloaking device on your golem. Your golem can use this device to cast invisibility on itself, without expending a spell slot or requiring components.

Once the golem casts the spell in this way, it can’t do so again until you finish a short or long rest.

\Subheading{Expanded Frame}
 
You enlarge your golem, increasing its size by one category, up to a maximum of Large. If your golem is Large after taking this upgrade, it gains advantage on Strength checks and Strength saving throws, and its hit point maximum increases to (5 + (the golem's Constitution modifier +5) * your inventor level).

During the process, you can determine if your golem will have the appropriate anatomy to serve as a mount. You can select this upgrade multiple times, but you can’t increase the golem’s size above Large using this upgrade.

\Subheading{Iron Fortress Prerequisite: golem of Medium or larger size with a Constitution of 18 or higher}
 
You extend your golem’s shielding and stationary stability. Your golem now counts as three-quarters cover for creatures riding it or for creatures within 5 feet of it (as long as it is between the target and the attacker). Additionally, while your golem is on the ground, it can’t be moved against its will.

\Subheading{Reciprocity Programming}
 
If you use your bonus action to take the Help action to grant your golem advantage on an attack against a creature, your golem can take the Help action as a bonus action on its next turn to grant you advantage on an attack against a creature.

\Subheading{Shielding Bond}
 
After studying the arcane fundamentals of Shield Golems, you have gained insight on how they shield their controllers, and can implement it in your own golem. Your golem gains the ability to cast warding bond without expending a spell slot. When cast in this method, it doesn’t require material components.

Once the golem casts warding bond in this way, it can’t use this feature to cast it again until it finishes a short or long rest.

\subsection{9th-Level Upgrades}

\Subheading{Mark of Sapience (Incompatible: Launcher Golem Chassis)}
 
You have attained the understanding of magic to craft a Mark of Sapience on the forehead of your Mechanical Golem, turning it into a Golem Companion. Its mental ability scores increase, granting it an Intelligence of 10, a Wisdom of 10, and a Charisma of 8. This allows it to speak, remember things, and follow more complex commands without direct input. Additionallly, it gains proficiency in Intelligence and Wisdom saving throws.

\begin{CalloutBox}{Roleplaying an Ironwrought Companion}

If you choose the Mark of Sapience upgrade, your Ironwrought Golem becomes a sapient companion, capable of learning, thinking, and having opinions. Consider how this may impact your interactions.

Fundamentally, a Ironwrought Companion is still entirely loyal to its creator, but it can develop a personality and thoughts of its own.
\end{CalloutBox}

\Subheading{Overdrive}
 
You build in a special mode that allows your golem to push beyond its limitations. As an action, you can overcharge your golem with energy, granting it the effects of the haste spell for a number of rounds equal to your Intelligence modifier.

Once you overcharge your golem in this way, you can’t do so again until you finish a long rest.

\Subheading{Powered Charge}
 
You improve your golem’s charging capabilities, increasing its walking speed by 5 feet. If your golem has a walking speed of 40 feet or higher after taking this upgrade, it gains the Forceful Slam trait below. You can select this upgrade up to two times.
\begin{fiveitemize}
	\item \textbf{Forceful Slam.} If the golem moves at least 20 feet straight toward a target and then immediately hits it with a melee attack, the target must succeed on a Strength saving throw (DC 8 + the golem’s proficiency bonus + its Strength modifier) or be knocked prone.
\end{fiveitemize}

\subsection{11th-Level Upgrades}

\Subheading{Airborne Propulsion}
 
You add a new method of propulsion to your golem, such as intricate mechanical wings or propellers. It gains a flying speed of 30 feet.

\Subheading{Multiattack Protocol}
 
Your golem can attack twice, instead of once, whenever it takes the Attack action on its turn.

\Subheading{Thundering Stomp (Prerequisite: Golem of Large size)}
 
Your golem can leverage its increased size and magical nature, allowing it to replace any attack it makes as part of the Attack action with a crushing stomp of magical energy as it brings down its foot. When it does so, each creature within 5 feet of the golem, other than itself, must succeed on a Constitution saving throw against your spell save DC or take 1d6 thunder damage.

\Subheading{Transforming Golem}
 
You install clever, multifunctional components into your golem, allowing it to reduce its size. Your golem can collapse itself down, decreasing its size by one category—from Medium to Small, for example. While collapsed in this way, its speed is reduced by 10 feet. It takes 1 minute for your golem to collapse itself or revert to its normal size.


\subsection{15th-Level Upgrades}

\Subheading{Artificial Learning (Prerequisite: Mark of Sapience)}
 
Your Golem Companion begins to apply its abilities to learn new things, gaining one level in a class of your choice. Your Golem Companion gains all the 1st-level features of the chosen class. This doesn’t include hit points, Hit Dice, or class proficiencies. For example, choosing the Fighter class grants the golem the Fighting Style and Second Wind features, but nothing else.

You can select this upgrade multiple times, granting your golem another level of features from the same class or the 1st-level features from another class each time it is selected.

\Subheading{Brutal Armaments}
 
You increase the effectiveness of your golem’s natural weapons, allowing it to attack with lethal strikes. When your golem makes an attack, you can choose to subtract its proficiency from the attack roll. If you do so, and the attack hits, you can add double its proficiency bonus to the damage roll of the attack.

\Subheading{Shared Power}
 
You bind your golem’s power source to your own soul, allowing you to tap its power and it to tap your power. Your golem can use its action and expend one of your spell slots to cast a spell you know of the slot’s level or lower.

In addition, you can use your action to cause your golem to lose a number of hit points up to twice your inventor level, granting you a number of temporary hit points equal to that amount.

Finally, on your turn, you can forgo your action to grant your golem an additional action, or you can command your golem to forgo its action, granting yourself an additional action. That action can only be used to take the Attack (one weapon attack only), Dash, Disengage, Hide, or Use an Object action.

\Subheading{Titan Slayer}
 
You build into your golem one oversized weapon. This weapon can be a martial melee weapon or an improved version of any natural weapon the golem has. The weapon’s base damage dice are doubled—from 1d8 to 2d8, for example. The golem can only attack with this oversized weapon once per turn (making any additional attacks it can make using other weapons).

\SectionSplashPage[0.4\textheight]{Infusionsmith.png}
\begin{multicols}{2}
\section{Infusionsmith}
An infusionsmith is, in some ways, perhaps the most quintessential type of inventor. While other inventors may delve into mechanics and tinkering, an infusionsmith tinkers with magic itself.

These inventors are on the cutting edge of magical engineering, understanding the principal applications of magic. Infusionsmiths would have ground to stand on in calling a wizard an impulsive spellslinger, for they are the inventors that work their magic through careful and meticulous study, laying down magic they may not use for hours, or painstakingly crafting a long-lasting enchantment.

An infusionsmith can be a magical swordsman, a wandslinger, or a bookish supporting character with a bagful of tricks that never quite runs dry. Some will stand within a maelstrom of animated blades, while others will plink away with a crossbow that seems unerringly accurate, but their common theme is using their prowess of infusion to make themselves formidable.
\subsection{Infusionsmith's Proficiency}

When you choose this specialization at 1st level, you gain proficiency with calligrapher’s supplies and jeweler’s tools. Your knowledge of infusion magic also gives you a natural affinity for scribing spell scrolls. The time you must spend to create a magic spell scroll is halved.
\subsection{Infused Armament}

Starting at 1st level, you can infuse a powerful armament to take into battle. At the end of a long rest, select one of the infusions below. The infusion lasts until the end of your next long rest.
\end{multicols}
\newpage
\twocolumn

\Subheading{Animated Weapon}

You touch a melee weapon, causing it to spring to life. This Animated Weapon can be carried or stowed like a normal weapon, or you can ready it, causing it to float beside you. While an Animated Weapon is readied, you can make attacks with it as part the Attack action, sending it out to strike a target. All attacks made with your Animated Weapon during a turn must be made against the same target.

This special attack is a melee spell attack with which you are proficient. You can make this special attack against a target out to a range of 30 feet away from you. If the weapon has the heavy or special property, this range is halved to 15 feet; and if the weapon has the light property, this range is doubled to 60 feet. On a hit, the target takes damage equal to the weapon’s damage dice + your Intelligence modifier.

\Subheading{Blasting Rod}

You touch a nonmagical object—a wand-blank, stick, staff, or rod—turning it into a Blasting Rod and infusing it with the ability to cast a cantrip. Select one evocation cantrip from the wizard spell list that doesn’t require concentration. Thereafter, as an action, you can use the Blasting Rod to cast that cantrip.

Once per turn, when you deal damage to a creature or object with your Blasting Rod, you can add your Intelligence modifier to the damage dealt to that target.

\Subheading{Infused Weapon}

You touch a weapon, enchanting it. While this weapon is enchanted, you (and only you) have proficiency with it. This Infused Weapon can be wielded like a normal weapon, but it gains the following special property:

You can use your Intelligence modifier, instead of Strength or Dexterity, for the weapon’s attack and damage rolls. In addition, if the weapon has a single damage die, the size of that die increases by one (to a maximum of a d12). For example, if the Infused Weapon is a dagger, its damage die increases from a d4 to a d6.

\subsection{Spell Manual}

At 3rd level, you have a Spell Manual containing two 1st-level wizard spells of your choice. Your Spell Manual is a repository of any non-inventor spell you know.

You can’t cast spells from your Spell Manual, and they don’t count against your number of spells known, but whenever you would learn a new inventor spell, you can select a spell from your Spell Manual instead of the inventor spell list. That spell then becomes an inventor spell for you, and you can cast it as normal. You can choose to scribe any inventor spell you know into your Spell Manual, as if copying it from a written spell.

Learning Spells of 1st Level and Higher. Each time you gain an inventor level, you can add one wizard spell of your choice to your Spell Manual for free. The chosen spell must be of a level for which you have spell slots, as shown on the Inventor table.

On your adventures, you might find other spells that you can add to your Spell Manual. For each level of the spell, the process takes 2 hours and costs 50 gp.

\subsection{Infuse Magic}

Also at 3rd level, you can use your Spell Manual to infuse magic items, such as wands, storing magical power for later use. Over the course of 1 minute, you can perform a special ritual to cast a spell into the item, expending a spell slot as normal. The spell being stored has no immediate effect when cast in this way.

Thereafter, a creature holding the item that is aware there is magic infused in it can trigger the stored magic, casting the spell. The creature must have an Intelligence of 6 or higher to cast the spell in this way.

The spell uses your spell save DC, spell attack bonus, and spellcasting ability, but is in all other ways treated as if the creature holding it cast the spell. The spell infused in the item fades if you finish a long rest and the stored magic hasn’t been triggered.

\subsection{Extra Attack}

Beginning at 5th level, you can attack twice, instead of once, whenever you take the Attack action on your turn.

\subsection{Empowered Infusions}

By the time you reach 5th level, your Infused Armament grows more powerful. If you infuse or animate a weapon, you can infuse or animate one additional weapon at the same time (having two Animated Weapons, two Infused Weapons, or one of each). If you have multiple Animated Weapons, you can attack the same or different targets with them during a turn.

If you cast a spell of 1st level or higher using an item from Infuse Magic or an upgrade, or from a magic wand, you can cast the cantrip from your Blasting Rod as a bonus action.

\subsection{Infused Focus}

Starting at 14th level, you can anchor a powerful spell into an item. When you cast a spell that requires concentration, you can anchor it to an object you touch (such as a staff, wand, or weapon), causing the spell to not require your concentration. When you do so, the spell lasts a number of rounds equal to your Intelligence modifier, after which the spell ends.
\begin{CalloutBox}{Benefits of Flexibility?}
If you choose to have an Animated and Infused Weapon, you can attack twice with either, but will limit yourself to being able to only attack one target if you solely rely on the Animated Weapon. If you animate two weapons, you can split your attacks like normal using one to attack a different creature. Alternatively, you can animate a short-range heavy weapon and a long-range light weapon, giving you more options at different ranges but limiting your ability to target multiple targets at any optimal range.
\end{CalloutBox}

Once you use this feature, you can’t use it again until you finish a short or long rest.
\section{Infusionsmith Upgrades}
\subsection{Unrestricted Upgrades}

\Subheading{Animated Archer}
 
You master animation enchantments, allowing you to use your Animate Weapon feature on a ranged weapon and a quiver of ammunition. An animated ranged weapon hovers near you, and you can make a ranged spell attack with it against a target within the weapon’s range.

On a hit, the target takes damage equal to the weapon’s damage dice + your Intelligence modifier. The weapon requires ammunition, and it can carry up 30 pieces of ammunition at a time. You can reload the ammunition at any time as an action.

\Subheading{Arcane Armament}
 
You master armoring yourself with magical enchantments. You learn the mage armor spell. While under the effect of mage armor , you can add your Intelligence modifier to your AC, instead of your Dexterity modifier. You can only do this if you are not adding your Intelligence modifier to your AC from another source.

Additionally, you gain resistance to force damage.

\Subheading{Loyal Weapon}
 
You can cast the returning weapon* spell without expending a spell slot, but if you cast it on a second weapon with this upgrade, the spell immediately fades from any previous use of this upgrade.

\Subheading{Magical Wand of...}
 
You create a new wand that you can infuse with a spell of 1st level or higher that you have recorded in your Spell Manual. This wand doesn’t require attunement, but can only be used by you. The spell must be of a level that you can cast as an inventor (as of when you would choose this upgrade). This wand has 3 charges.

You can expend a charge to cast the chosen spell at its lowest level. The wand regains all expended charges when you finish a long rest.

You can select this upgrade multiple times, choosing a different spell and creating a new wand each time you select this upgrade.
\begin{CalloutBox}{Magic Wand Spell Level}
A Magical Wand created using an upgrade slot at level 3 can always only ever cast 1st-level spells. For example, at level 5 if you have two Magical Wand upgrades, you will have one that can cast 1st-level spells, and one that can cast 2nd-level spells. Even if you later swap the upgrade, the Wand is still an Upgrade from level 3, and selected as if you are level 3 and only able to cast 1st-level spells.
\end{CalloutBox}

\Subheading{Skilled Infusions}
 
You manage to make the magic of your Animated and Infused Weapons so potent that attacks made with them gain the benefits of a fighting style. Attacks made with one-handed melee weapons gain Dueling, attacks made with two-handed melee weapons gain Great Weapon Fighting, and attacks made with ranged weapons gain Archery.

\Subheading{Soul-Saving Bond}
 
You set up a special magical bond between you and another creature. When either of the bonded creatures fails an Intelligence, Wisdom, Charisma, or death saving throw, the other creature can make the same saving throw, replacing the original result with its own. If this ability is used on a death saving throw, the roll is automatically a 20.

Once this ability has been used to replace a roll, it can’t be used again until both bonded creatures have finished a short or long rest.

This bond can be set up with a different creature at the end of a long rest.

\Subheading{Warding Stone}
 
You learn how to weave a protective enchantment on an item. That item gains a pool of temporary hit points equal to your Inventor level. Any creature carrying the item gains any temporary hit points remaining in its pool, but these hit points are lost if that creature is no longer carrying the item.

The item’s pool of temporary hit points replenishes when you finish a long rest.

\Subheading{Worn Enchantment}
 
You can enchant an item you are wearing, such as a cloak or scarf, to animate and assist you with a task, be it climbing a wall, grappling an enemy, or picking a lock. You can expend a spell slot of 1st level or higher as an action to gain proficiency in one Strength- or Dexterity-based skill of your choice until you finish a long rest.

Additionally, you can choose to use up all the magic in the item to gain advantage on one ability check using the chosen skill. After you do so, you lose the proficiency granted by this upgrade.

\subsection{5th-Level Upgrades}

\Subheading{Animated Shield}
 
You apply your animation magic to a shield. At the end of a long rest, you can touch a shield, causing it to spring to life and protect you until you finish your next long rest. While it is protecting you, you gain the benefits of the shield as if you were wielding it. As usual, you can benefit from only one shield at a time.

While you have an animated shield, you can also send it to defend your allies. When a creature you can see attacks a target other than you that is within 30 feet of you, you can use your reaction to impose disadvantage on the attack roll. When you do so, you lose the benefits of the animated shield until the start of your next turn.

\Subheading{Arcane Ammunition}
 
You infuse a ranged weapon with special, arcane magic. The weapon no longer requires ammunition to fire, as it creates magical ammunition when you fire it, and whenever you hit a target with this magical ammunition, the attack deals force damage, instead of the weapon’s normal damage type.

Additionally, if the weapon had the loading property, it no longer has that property.

\Subheading{Deflecting Weapon}
 
While you have an animated melee weapon, you can use your reaction when a creature makes a weapon attack against you to attempt to block it, rolling a d8 and adding the number rolled to your AC for that attack. You make this choice after the creature makes its roll, but before the GM says whether the attack hits or misses. If you attempt to block a melee weapon attack in this way and the attack misses, you can immediately make one weapon attack against the attacker as part of the same reaction.
When you use this ability, if you take the Attack action on your next turn, you can only make one weapon attack as part of that action. This attack can’t be made with the animated weapon you used to attempt to block the attack.

\Subheading{Lesser Ring of Protection}
 
You make a prototype ring of protection . Any creature wearing this ring gains a +1 bonus to its AC.

\Subheading{Malicious Infusion}
 
You learn the heat metal spell.

Additionally, as a reaction to being hit by a metal weapon, you can cast heat metal targeting the weapon that struck you, expending a spell slot as normal.

\Subheading{Ring of Reaction}
 
You make a minor magic ring. Any creature wearing this ring can add your Intelligence modifier to its Dexterity saving throws and initiative rolls.

\Subheading{Translocation Binding}
 
When you make an attack with an Animated or Infused melee weapon, you can use your bonus action and expend a spell slot of 1st level or higher to teleport to an unoccupied space within 5 feet of the weapon before it returns to you. You can use this ability even if your attack misses the target.

\Subheading{Weapon Enchantment Expertise}
 
When you cast arcane weapon* , magic weapon , or vorpal weapon* , the targeted weapon deals an extra 1d4 force damage, and when you cast prismatic weapon* , the extra damage it deals increases by 1d4.

Additionally, you have advantage on Constitution saving throws made to maintain concentration on arcane weapon*, prismatic weapon* , magic weapon , and vorpal weapon* .

\Subheading{Weapon Enchantment Resonance}
 
When you cast arcane weapon* , prismatic weapon* , magic weapon , or vorpal weapon* on one of your Animated or Infused weapons, all of your Animated and Infused weapons gain the benefits of that spell. Additionally, you can target any weapon with arcane weapon*, prismatic weapon* , or magic weapon , even if the weapon is already magical.

\subsection{9th-Level Upgrades}

\Subheading{Detonate Armament (Prerequisite: Arcane Armament)}
 
When you take damage while under the effects of mage armor , you can end the effects of that spell as a reaction to cast turbulent warp* without expending a spell slot.

Once you cast the spell in this way, you can’t do so again until you finish a short or long rest.

\Subheading{Dimensional Pockets}
 
You enchant a pocket on your gear to contain far more than it would appear. That pocket can hold up to 50 pounds, and you can put any object in it that can fit through a 6-inch-diameter opening. The pocket otherwise behaves like a bag of holding.

\Subheading{Invisibility Cloak}
 
You make an invisibility cloak. While wearing this cloak, a creature can cast invisibility on itself, without expending a spell slot or material components. When cast in this way, the spell doesn’t require the caster’s concentration. The spell ends if the caster is no longer wearing the cloak or if they choose to end the spell early (no action required).

Once a creature uses the cloak to cast this spell, it can’t be used in this way again until you finish a short or long rest.

\Subheading{Prepared Enchantment}
 
You can cast a limited version of contingency , without expending a spell slot or material components, as you bestow magical enchantments onto your gear. When you cast the spell in this way, it ends if you finish a long rest before the contingent spell is triggered, and the contingent spell must be a spell of 2nd level or lower.

Once you cast contingency in this way, you can’t do so again until you finish a long rest.

\Subheading{Spell-Trapping Ring}
 
At the end of a long rest, you can set a powerful magic into a nonmagical ring you touch. You can use this ring to cast counterspell without expending a spell slot. When you cast counterspell in this way and it succeeds, the spell you countered is stored in the ring. You can cast the stored spell from the ring without expending a spell slot, but that spell fades if you don’t cast it before you finish a long rest.

Once you cast counterspell to store a spell in this ring, you can’t do so again until you finish a long rest.

\subsection{11th-Level Upgrades}

\Subheading{Empower Weapon}
 
When you hit a creature with a weapon, you can use your bonus action to infuse your strike with arcane energy, dealing an extra 2d6 force damage to the target. When you use this bonus action, you can also choose to expend a spell slot of 1st level or higher to further increase this extra damage. When you do so, the target takes an extra 1d6 force damage, plus another 1d6 per level of the spell slot.

\Subheading{Enchanted Broom}
 
You can enchant a broom (or broom-like object), turning it into a broom of flying . You set the command word for the broom, and it only obeys you.

\Subheading{Life Infusion}
 
You learn a potent magical infusion that suffuses a creature with life energy. You can cast regenerate , without expending a spell slot or material components.

Once you cast the spell in this way, you can’t do so again until you finish a long rest.

\Subheading{Magical Rod of...}
 
You create a new rod that you can infuse with a spell of 5th level or higher that you have recorded in your Spell Manual. This rod doesn’t require attunement, but can only be used by you. The spell must be of a level no higher than half your inventor level (as of when you would choose this upgrade), rounded down.

You can use the rod to cast the chosen spell at its lowest level. Once you cast the spell in this way, you can’t do so again until you finish a long rest. You can select this upgrade multiple times, choosing a different spell and creating a new rod each time you select this upgrade.

\begin{CalloutBox}{How to get high-level spells?}
Astute players will note that when you add spells to your Spell Manual through leveling, the spell must be of a level for which you have spell slots, but when you gain access to the Magical Rod upgrade, you don’t yet have 5th-level spell slots. Spells for Magical Rods will primarily have to be found in the wild, in the form of scrolls or spellbooks, and copied into your Spell Manual that way.
\end{CalloutBox}

\Subheading{Mixed Technique}
 
When you use your action to cast the cantrip from your Blasting Rod, you can make one weapon attack with an Animated or Infused Weapon as a bonus action.

\subsection{15th-Level Upgrades}

\Subheading{Advanced Object Animation}
 
When you cast the animate objects spell, the animated objects use your spell attack bonus for their attack rolls, and they gain a bonus to their damage rolls based on their size: +1 (Tiny or Small), +2 (Medium), +4 (Large), or +8 (Huge).

\Subheading{Third Animated Weapon}
 
Your mastery of weapon animation expands to a greater breadth of control. If you have two Animated Weapons, you can animate a third weapon with your Animate Weapon feature. When you take the Attack action and make an attack with your other two Animated Weapons, you can make an attack with your third Animated Weapon as part of the same action.

\begin{CalloutBox}{High-Level Rods \& Wands}
While it may appear that infusionsmiths lack high-level upgrades, their Magical Rod and Wand upgrades provide a plethora of very valuable, high-level options.
\end{CalloutBox}

\SectionSplashPage[0.4\textheight]{Potionsmith.png}
\begin{multicols}{2}
\section{Potionsmith}
A potionsmith is an inventor who has pursued the secrets of alchemy. While many a village has an apothecary grinding odd herbs and roots into potent (or not so potent) concoctions, the careful process of mixing and brewing is just one way to achieve results. Through the use of the intricate secrets of the craft and direct infusions using magical rituals, a potionsmith can come up with explosive results in the blink of an eye... sometimes literally.

A potionsmith can be a scholar who has delved into the knowledge the world has to offer or an explorer that has unlocked the secrets of the wilderness. Their knowledge could come from being friends of the fey or from unfettered access to the royal library. As such, potionsmith’s can be good or evil, lawful or chaotic.
\subsection{Potionsmith's Proficiency}

When you choose this specialization at 1st level, you gain proficiency with blowguns, alchemist’s supplies, and the herbalism kit.

Your knowledge of alchemy gives you a natural affinity for brewing potions. Creating a potion through normal crafting takes you only half the time it would normally take.

\subsection{Alchemical Reagents Pouch}

At 1st level, you’ve acquired a pouch of useful basic reagents, much as a wizard might carry a component pouch. You’ve tucked away things that will come in handy, and can retrieve them as part of using an ability that might require them. As long as you have this pouch on you, you can use the potionsmith’s features. This pouch is considered to be alchemist’s supplies for the purposes of crafting.

If you lose your reagent pouch, you can spend 50 gp to reacquire the various stocks you need, or spend 1 full day gathering them from a natural environment without expense.
\end{multicols}
\newpage
\twocolumn
\begin{CalloutBox}{Alchemists and Potions}
Immediately on making a potionsmith, one may think of the magic potions that exist in most settings—from the simple potion of healing to the legendary potion of storm giant strength. However, someone does not need to be a potionsmith to make these, as crafting rules for them are open to all classes.

A potionsmith with their proficiency will certainly have the edge in potion crafting, but these potions are not fundamental to the class.
\end{CalloutBox}
\subsection{Instant Reactions}

At 1st level, you know how to get instant reactions to occur without the niceties of grinding, simmering, and brewing required for more refined concoctions.

Moreover, you know enough ways to do these processes that with just a few standard supplies, you can get results out of a wide range of things you can gather in almost any locale and a pinch from your reagent pouch.

Pick three instant reactions from the following: Alchemical Acid, Alchemical Fire, Fortifying Fumes, Healing Draught, and Poisonous Gas. These instant reactions don’t count against your number of upgrades known.

\subsection{Alchemical Infusions}

At 3rd level, you have discovered the secrets to drastically bypass the alchemical brewing process, crafting potent magical effects through the direct infusion of the potion with your own magic. When you finish a short or long rest, you can create any number of potions, choosing one of the following spells for each potion and expending a spell slot as if you were casting the spell. You must provide an empty vial for each potion you create in this way.

\vspace{1pt}
\begin{SmallTableBox}
\small % optional, helps in two-column mode
\renewcommand{\arraystretch}{1.2}
\rowcolors{2}{AltRow}{Parchment}
\begin{tabularx}{\linewidth}{@{}>{\raggedright\arraybackslash}p{2.5cm}X@{}}
\rowcolor{TableHeader}
\textbf{Spell Level} & \textbf{Infusion Spells} \\
1st  & \textit{cure wounds, fog cloud, grease, heroism} \\
2nd  & \textit{barkskin, lesser restoration, shatter, web} \\
3rd  & \textit{blink, haste, stinking cloud, water breathing} \\
4th  & \textit{confusion, stoneskin} \\
5th  & \textit{cloudkill} \\
\end{tabularx}
\end{SmallTableBox}
\vspace{1pt}

When you do so, the spell doesn’t take effect immediately, but is instead infused into the potion.

If the spell grants an effect or restores hit points, a creature can consume it or administer it to another creature as an action. If a spell targets an area of effect, you can accurately throw the vial up to 30 feet, dealing no inherent damage, but breaking the vial on impact. The area of effect takes place when the vial is broken, with the effect centered on where the vial breaks. If the spell targets a creature, you can treat the vial as a ranged weapon with which you are proficient, and it has the thrown (20/60) property. You replace the spell attack with making a weapon attack made with the vial.

If the spell has a non-instantaneous effect that normally requires concentration, it doesn’t require your concentration when it takes effect in this way, but its duration is shortened to a number of rounds equal to your Intelligence modifier (minimum 1). A spell that doesn’t require concentration lasts its normal duration. An infused potion loses its potency if it isn’t used by the end of your next long rest.

You can gain additional spells for your Alchemical Infusions through your potionsmith upgrades.
\begin{CalloutBox}{Creativity and Catapults}
While you can only accurately throw the vial containing an infused potion up to 30 feet, as per the feature, the trigger merely specifies that the effect takes place where it breaks.

If you can find another way to deliver the vial, such as the launch object* spell (or a similar spell) on the inventor spell list—or, well, an actual catapult—you can apply these methods as well.
\end{CalloutBox}

\subsection{Practiced Quaff}

Additionally at 3rd level, you gain the ability to consume potions and infused potions as a bonus action.

\subsection{Empowered Alchemy}

Starting at 5th level, whenever you deal damage, grant temporary hit points, or restore hit points to a creature with an instant reaction or infused potion on your turn, you can add your Intelligence modifier to the amount of damage dealt or the number of hit points granted or restored.

\subsection{Infusion Expertise}

Starting at 14th level, when you create an infused potion at end of a short or long rest, the first infused potion you create doesn’t require a spell slot to infuse, and you can choose a spell for which you would not otherwise have a high enough level spell slot when making this infusion.

\section{Potionsmith Upgrades}
\subsection{Unrestricted Upgrades}

\Subheading{Alchemical Acid (Instant Reaction)}
 
As an action, you can produce an instant reaction, throwing this combination of reagents at a point within 20 feet of you, which causes a caustic acid to splatter the area within a 5-footradius sphere centered on that point. Each creature in that area must succeed on a Dexterity saving throw against your spell save DC or take 2d4 acid damage. Damage from this acid deals double damage to objects and structures.

The damage increases by 2d4 when you reach 5th level (4d4), 11th level (6d4), and 17th level (8d4).

Alternatively, you can use a bonus action to prepare this instant reaction as a ranged weapon with which you are proficient. It lasts until the end of your turn, and it has the finesse and thrown (20/60) properties. On a hit, it deals damage equal to failing the saving throw against this effect. The damage dealt by this attack doesn’t include your Strength or Dexterity modifier.

\Subheading{Alchemical Fire (Instant Reaction)}
 
As an action, you can produce an instant reaction, tossing a quick combination of reagents at a point within 20 feet of you, which causes searing fire to flare up in a 5-foot-radius sphere centered on that point. Each creature in the area must succeed on a Dexterity saving throw against your spell save DC or take 1d8 fire damage.

This damage increases by 1d8 when you reach 5th level (2d8), 11th level (3d8), and 17th level (4d8).

Alternatively, you can use a bonus action to prepare this instant reaction as a ranged weapon with which you are proficient. It lasts until the end of your turn, and it has the finesse and thrown (20/60) properties. On a hit, it deals damage equal to failing the saving throw against this effect. The damage dealt by this attack doesn’t include your Strength or Dexterity modifier.

\Subheading{Delivery Mechanism}
 
You modify the stability of your reagents and develop a better delivery mechanism. You can target a point within 40 feet of you for your instant reactions (including ones gained from upgrades, such as Explosive Reaction) and infused potions that target a point. This additional precision enables you to better target the effects, allowing each creature of your choice in the area to automatically succeed on its Dexterity saving throw against the effects of your instant reactions and infused potions.

Additionally, when making an attack roll with an instant reaction or infused potion, you can use your Intelligence modifier, instead of your Dexterity modifier, for the attack roll.

\Subheading{Dragon Draught (Instant Reaction; upgrade only)}
 
As an action, you can quaff a vial of highly reactive liquid, allowing you to immediately belch a blast of devastating elemental energy. Choose one of the following damage types: acid, cold, fire, or lightning.

If you choose cold or fire damage, the area of effect is a 10-foot cone; if you choose acid or lightning damage, the area of effect is a 15-foot line. Each creature in the area must succeed on a Dexterity saving throw against your spell save DC or take 1d8 damage of the chosen type.

The damage increases by 1d8 when you reach 5th level (2d8), 11th level (3d8), and 17th level (4d8).

\Subheading{Explosive Reaction (Instant Reaction; upgrade only)}
 
As an action, you can produce an instant reaction, throwing a devastating combination of reagents at a point within 20 feet of you, which causes a minor explosion. Each creature within 10 feet of that point must succeed on a Constitution saving throw against your spell save DC or take 1d10 thunder damage from the shockwave of the explosion.

The damage increases by 1d10 when you reach 5th level (2d10), 11th level (3d10), and 17th level (4d10).

Alternatively, you can use a bonus action to prepare this instant reaction as a ranged weapon with which you are proficient. It lasts until the end of your turn, and it has the finesse and thrown (20/60) properties. On a hit, it deals damage equal to failing the saving throw against this effect. The damage dealt by this attack doesn’t include your Strength or Dexterity modifier.

\Subheading{Flaming Grease}
 
You tweak your grease formula to be flammable, as grease should be. When you cast grease or use it to make an infused potion, the effect becomes flammable. If the grease takes fire damage, all creatures in the spell’s area take 2d4 fire damage. The grease burns for a number of rounds equal to your Intelligence modifier (minimum 1).

Additionally, when you cast the grease spell or use it as an infused potion, you can instead use it to coat a weapon. This coating lasts for 1 hour or until it takes fire damage, igniting the grease. Once ignited, that weapon deals an extra 1d4 fire damage on a hit for the next minute.

\Subheading{Fortifying Fumes Reaction (Instant Reaction)}
 
As an action, you can produce an instant reaction, throwing a fortifying stimulus of reagents at a point within 20 feet of you, which causes fumes to erupt. Each creature within 10 feet of that point can choose to hold its breath and not inhale the fumes, but creatures that do inhale them gain 1d4 temporary hit points, deal an extra 1d4 damage on their next melee weapon attack, and have advantage on their next Constitution saving throw. Any remaining benefits from this effect fade at the end of your next turn.

Both the temporary hit points and the extra damage increase by 1d4 when you reach 5th level (2d4), 11th level (3d4), and 17th level (4d4).

\Subheading{Frostbloom Reaction (Instant Reaction; upgrade only)}
 
As an action, you can produce an instant reaction, throwing a combination of devastating, endothermically reactive reagents at a point within 20 feet of you, which causes a bloom of ice to erupt. The area within 5 feet of that point becomes difficult terrain until the end of your next turn.

Each creature in the area must succeed on a Dexterity saving throw against your spell save DC or be caught by the ice, taking 1d6 cold damage. If a creature that fails its saving throw is entirely in the area of effect, it becomes restrained until the end of its next turn. A restrained creature can use its action to make a Strength saving throw against your spell save DC, ending the effect on itself on a success.

The damage increases by 1d6 when you reach 5th level (2d6), 11th level (3d6), and 17th level (4d6).

\Subheading{Healing Draught (Instant Reaction)}
 
As a bonus action, you can produce an instant reaction that provides potent magical healing. A creature can use its action to consume this draught or administer it to another creature within 5 feet of it before the start of your next turn, causing the target to regain 1d8 hit points. If you have the Practiced Quaff feature, you can consume a Healing Draught as part of the bonus action used to make it.

A creature can benefit from a number of these draughts up to its Constitution modifier (minimum of 1), after which they provide no benefit to that creature until it finishes a long rest. A Healing Draught that isn’t consumed before the start of your next turn loses its potency.

The number of hit points the draught restores increases by 1d8 when you reach 5th level (2d8), 11th level (3d8), and 17th level (4d8).

\Subheading{Homunculus Familiar}
 
You learn the alchemical process for creating a homunculus minion. You can cast find familiar without expending a spell slot or material components other than your alchemical reagents and blood. This familiar can appear in any Tiny shape you want, though it uses the statistics of a creature that can normally be chosen by the spell, and it is a construct or monstrosity, rather than one of the types listed in the spell.

\Subheading{Inoculations}
 
You gain resistance to poison damage.

Additionally, at the end of a long rest, you can choose up to five creatures (including yourself) to become inoculated against poisonous effects you can produce that require a Constitution saving throw (such as the Poisonous Gas instant reaction or the cloudkill alchemical infusion). You and each of the chosen creatures automatically succeed on saving throws against these effects until the end of your next long rest.

\Subheading{Long Acting}
 
If the spell you infuse into an infused potion would normally require concentration, its duration is now a number of rounds equal to your Intelligence modifier + your proficiency bonus.

\Subheading{Persistent Reactions}
 
Your instant reactions that have an area of effect persist in that area until the start of your next turn. While they persist, each creature that enters the area for the first time on a turn or ends its turn there must repeat the saving throw against the effect. You can choose to make the effect not persist in this way when you take the action to produce the instant reaction.

\Subheading{Poisoner’s Proficiency}
 
You delve into the darkest, most mysterious depths of herbal lore, learning the potent secrets of poison. You gain proficiency with poisoner’s kits; if you already have this proficiency, your proficiency bonus is doubled for any ability check you make using a poisoner’s kit. Additionally, during a long rest, you can create one of the following poisons of your choice:

\begin{fiveitemize}
	\item \textbf{Contact Poison.} You can apply this poison to a melee weapon or up to ten pieces of ammunition. That weapon or ammunition deals an extra 1d4 poison damage on a hit; the weapon can deal this poison damage 10 times before the poison wears off.
	\item \textbf{Ingested Poison.} This poison is a simple, flavorless powder. If a creature consumes a full dose of the poison, it must make a Constitution saving throw with disadvantage against your spell save DC. After 1 minute, if the creature failed its saving throw, it takes 1d10 poison damage per level you have in this class, and it becomes poisoned until it finishes a long rest.
	\item \textbf{Inhaled Poison.} You can use this dose of poison to increase the size and potency of your Poisonous Gas instant reaction. When you do so, the fumes erupt in a 20-footradius sphere and the effect’s damage dice are doubled.
\end{fiveitemize}

Any unused poison loses its potency and no longer has any effect after you finish your next long rest.

\Subheading{Poisonous Gas (Instant Reaction)}
 
As an action, you can produce an instant reaction, tossing a quick combination of reagents at a point within 20 feet of you, which causes noxious fumes to erupt in a 10-foot-radius sphere centered on that point. Each creature in that area must succeed on a Constitution saving throw against your spell save DC or take 1d4 poison damage and become poisoned until the end of its next turn.
The damage increases by 1d4 when you reach 5th level (2d4), 11th level (3d4), and 17th level (4d4).

\Subheading{Secrets of Acid}
 
You learn the secrets of infusing acid into your infused potions. You add the following options to your list of available spells for your Alchemical Infusions:

\vspace{1pt}
\begin{SmallTableBox}
\small % optional, helps in two-column mode
\renewcommand{\arraystretch}{1.2}
\rowcolors{2}{AltRow}{Parchment}
\begin{tabularx}{\linewidth}{@{}>{\raggedright\arraybackslash}p{2.5cm}X@{}}
\rowcolor{TableHeader}
\textbf{Spell Level} & \textbf{Infusion Spells} \\
1st  & \textit{summon ooze*} \\
2nd  & \textit{acid arrow} \\
3rd  & \textit{erode*} \\
\end{tabularx}
\end{SmallTableBox}
\vspace{1pt}
Additionally, you can create one infused potion of summon ooze* at the end of a long rest without expending a spell slot.

\Subheading{Secrets of Fire}
 
You learn the secrets of infusing fire into your infused potions. You add the following options to your list of available spells for your Alchemical Infusions:

\vspace{1pt}
\begin{SmallTableBox}
\small % optional, helps in two-column mode
\renewcommand{\arraystretch}{1.2}
\rowcolors{2}{AltRow}{Parchment}
\begin{tabularx}{\linewidth}{@{}>{\raggedright\arraybackslash}p{2.5cm}X@{}}
\rowcolor{TableHeader}
\textbf{Spell Level} & \textbf{Infusion Spells} \\
1st  & \textit{faerie fire} \\
2nd  & \textit{unstable explosion*} \\
3rd  & \textit{fireball} \\
\end{tabularx}
\end{SmallTableBox}
\vspace{1pt}
Additionally, you can create one infused potion of faerie fire at the end of a long rest without expending a spell slot.

\Subheading{Secrets of Flight}
 
You learn the secrets of infusing the ability to escape the shackles of gravity into your infused potions. You add the following options to your list of available spells for your Alchemical Infusions:

\vspace{1pt}
\begin{SmallTableBox}
\small % optional, helps in two-column mode
\renewcommand{\arraystretch}{1.2}
\rowcolors{2}{AltRow}{Parchment}
\begin{tabularx}{\linewidth}{@{}>{\raggedright\arraybackslash}p{2.5cm}X@{}}
\rowcolor{TableHeader}
\textbf{Spell Level} & \textbf{Infusion Spells} \\
1st  & \textit{feather fall} \\
2nd  & \textit{levitate} \\
3rd  & \textit{fly} \\
\end{tabularx}
\end{SmallTableBox}
\vspace{1pt}
Additionally, you can create one infused potion of feather fall at the end of a long rest without expending a spell slot.

\Subheading{Secrets of Frost}
 
You learn the secrets of infusing frost into your infused potions. You add the following options to your list of available spells for your Alchemical Infusions:

\vspace{1pt}
\begin{SmallTableBox}
\small % optional, helps in two-column mode
\renewcommand{\arraystretch}{1.2}
\rowcolors{2}{AltRow}{Parchment}
\begin{tabularx}{\linewidth}{@{}>{\raggedright\arraybackslash}p{2.5cm}X@{}}
\rowcolor{TableHeader}
\textbf{Spell Level} & \textbf{Infusion Spells} \\
1st  & \textit{freezing shell*} \\
2nd  & \textit{cold snap*} \\
3rd  & \textit{flash freeze*} \\
\end{tabularx}
\end{SmallTableBox}
\vspace{1pt}
Additionally, you can create one infused potion of freezing shell* at the end of a long rest without expending a spell slot.

\begin{CalloutBox}{Secrets of...}
Almost any set of three spells (a 1st-level, 2nd-level, and 3rd level spell) that are either short-term buffs or deal area of effect damage can work, as long as there is a thematic connection between the spells.

Consult with your GM about additional options if you don’t see what you want present here.
\end{CalloutBox}

\Subheading{Weapon Coating}

As a bonus action, you can apply one of your instant reactions to a melee weapon or piece of ammunition. The next time you hit a target with that weapon or ammunition before the end of your next turn, the target is subjected to the effects of the applied instant reaction. The creature automatically takes the damage or healing associated with the instant reaction, but makes a saving throw as normal against any additional effects.

\textbf{The Implications}

There is an instant reaction that restores hit points, rather than deals damage. This can be applied via the Weapon Coating upgrade as well, though the weapon will still deal its damage on a hit. Perhaps your allies will forgive a blowgun dart coated in a healing draught...if you can hit with the attack.

\subsection{5th-Level Upgrades}

\Subheading{Adrenaline Serum}

You learn to create a potent serum. As a bonus action, you can consume a dose of this serum, increasing your Strength and Dexterity modifiers by an amount equal to your Intelligence modifier (up to a maximum of +5) for a number of rounds equal to your Constitution modifier (minimum of 1), after which all effects of the serum fade at the start of your next turn. While you are under the effect of this serum, you also gain the benefits of the heroism spell. When these effects end, your speed is halved and you can’t benefit from your Adrenaline Serum again until the start of your next turn.

\Subheading{Explosive Powder (Prerequisite: Explosive Reaction)}
 
You’ve managed to stabilize and refine your Explosive Reaction. You can prepare a number of Explosive Reactions up to your Intelligence modifier (a minimum of 1). These Explosive Reactions don’t detonate instantly, but rather last for 1 minute and will only detonate when exposed to fire. A creature can be affected by the damage from no more than two Explosive Reactions simultaneously, but structures take the full damage of all Explosive Reactions in their area.

\Subheading{Resistance Potion}
 
You create a formula for making a simplified temporary resistance potion. During a long rest, you can create a potion that grants the drinker resistance to acid, cold, fire, lighting, poison, or thunder damage (selected during the long rest) for 10 minutes when consumed. If not consumed, it loses potency when you finish your next long rest.

\subsection{9th-Level Upgrades}

\Subheading{Aroma Therapies}
 
You expand your alchemical knowledge to be able to produce incense and simmering reagents that grant benefits to those that inhale their fumes. A creature that spends a long rest surrounded by these fumes regains an additional 2d4 Hit Dice, its exhaustion level is reduced by 1d4 (instead of 1), and it is cured of any nonmagical diseases it’s suffering from.

\Subheading{Infusion Stone}
 
You use the secrets of alchemy to create an Infusion Stone. You can use this stone in the process of infusing potions in place of expending a spell slot level. The spell you infuse in this way must be of a level for which you have spell slots.
Once you use the Infusion Stone, you can’t use it again until you finish a long rest.

\Subheading{Mana Potion}
 
During a short rest, you can create a mana potion. A mana potion loses its potency if it isn’t consumed within 1 hour. As an action, a creature can consume a mana potion to regain one expended spell slot of its choice that is 3rd level or lower.

\Subheading{Potent Reactions}
 
You refine your instant reactions, increasing their potency. The die you roll to determine the damage or healing effect of your instant reactions is increased by one size: a d4 becomes a d6, a d6 becomes a d8, a d8 becomes a d10, and a d10 becomes a d12.

\Subheading{Rocketry (Prerequisite: Explosive Reaction and Delivery Mechanism)}
 
Combining your understanding of explosives with advances in delivery and stabilization, you produce the logical extreme: powering a rocket-based payload with your devastating knowledge.

You can choose any instant reaction you know as the payload for your rocket. Rockets must be prepared ahead of time, and you can prepare a number of them up to your Intelligence modifier (minimum of 1) at the end of a short or long rest. When you finish a rest, any unused rockets lose their potency and must be remade due to their volatile components.

A rocket targets a point you can see within 500 feet of you, but the DC of the instant reaction’s saving throw is reduced by 2 for every 100 feet it travels. When the rocket reaches its destination, the instant reaction takes effect, centered on that point.

Alternatively, you can load an alternate payload into your rockets, which can be an object that weighs no more than 1 pound, replacing the effect of the instant reaction. This alternate payload can be salvaged during a rest.

\subsection{11th-Level Upgrades}

\Subheading{Field Infusion}
 
You can create an infused potion as an action.

Once you create an infused potion in this way, you can’t do so again until you finish a short or long rest.

\Subheading{Mutation Mixture (Prerequisite: 13th-level Inventor)}
 
You concoct an infused potion that can rapidly warp a creature’s body, temporarily mutating it in extreme ways. You add the polymorph spell to your list of available spells for your Alchemical Infusions. When a creature consumes this infused potion, the spell has a duration of 1 hour, but the creature that consumed it can end the effect early as an action.

Once a creature consumes this infused potion, it can’t be affected by the polymorph spell in this way again until it finishes a long rest.

\Subheading{Panacea}
 
When you create a Healing Draught, you can add a more potent concoction to it. A creature that consumes this potent draught can treat any dice rolled to determine the hit points it regains as having rolled their maximum value, and the creature gains the benefit of a greater restoration spell.

Once you create a Healing Draught in this way, you can’t do so again until you finish a long rest.

\Subheading{Perfect Reaction}
 
In a moment of perfect focus, you can create a flawless instant reaction. When you deal damage, grant temporary hit points, or restore hit points to a creature with this instant reaction, you can use the maximum possible result instead of rolling.

Once you create a Perfect Reaction, you can’t do so again until you finish a short or long rest.

\subsection{15th-Level Upgrades}

\Subheading{Adrenaline Rush (Prerequisite: Adrenaline Serum)}
 
While under the effects of your Adrenaline Serum, you can attack twice, instead of once, whenever you take the Attack action on your turn.

\Subheading{Mad Alchemy}
 
Whenever you create an infused potion, you can infuse two spells into it, instead of one, which activate simultaneously.

\Subheading{Philosopher’s Stone}
 
You create a Philosopher’s Stone, allowing you to recreate wonders of alchemy. As long as you have a supply of non-gold metal, you can create up to 1 pound of gold each day (50 gp worth). Additionally, the Philosopher’s Stone can be used in place of a diamond worth up to 500 gp for a spell’s material component; it isn’t consumed when used in this way, but it becomes inert for 24 hours, providing no benefits.

You can brew a special potion using your Philosopher’s Stone called an Elixir of Life. Brewing this potion takes 8 hours and requires crushing a diamond worth at least 2,000 gp into it. An Elixir of Life causes a creature that drinks it to cease aging for 4d4 years. A creature that drinks this elixir also gains the effect of a death ward spell, which lasts until it is triggered, instead of 8 hours.

A more potent Elixir of Life can be created, increasing the number of years that the creature ceases aging by 1d4 for each additional diamond spent.

\SectionSplashPage[0.4\textheight]{Thundersmith.png}
\begin{multicols}{2}
\section{Thundersmith}
A Thundersmith is an inventor who harnesses the primal force of elemental power, channeling that power into their great creation: a weapon of unmatched devastation. Spectacular and terrible, these weapons bring fear to their foes and awe to their allies.

Why a Thundersmith bends their mind to the task of making such a weapon are as varied as the creations themselves. Some are coldly analytical about the destruction it causes seeking to continually improve it, tweaking it for ever more optimized destruction, while others view it merely as a tool, a means to an end; others still revel in the crash of thunder that heralds the terrifying force of their weapon.

Many look at these weapons as the dawn of a new age, in truth wielding them is a tricky and arcane art; as complex as its creation it is only truly understood and mastered by the one who forged the device—each weapon a unique piece of devastating art.
\subsection{Thundersmith's Proficiency}

When you choose this specialization at 1st level, you gain proficiency with tinker’s tools and smith’s tools.

If your weapon requires ammunition, you gain the knowledge of how to forge and create it with smith’s tools during a long rest. You can create up to 50 rounds of ammunition during a long rest, with materials costing 1 gp per 10 rounds.
\end{multicols}
\newpage
\twocolumn
% Overlay the singed background *before* placing Stormcharged.png
\begin{tikzpicture}[remember picture, overlay]
  \node[
	opacity=0.75,
    anchor=south east,
    inner sep=0pt
  ] at ([xshift=0cm, yshift=0cm]current page.south east)
  {\includegraphics[width=\paperwidth,height=\paperheight]{Singed_Background.png}};
\end{tikzpicture}

\subsection{Stormforged Weapon}

Starting at 1st level, you harness the elemental power of thundering storms to create a powerful weapon. This weapon requires attunement, only you can attune to it, you are proficient with it while attuned, and you can only be attuned to one Stormforged Weapon at a time. If you have multiple Stormforged Weapons, you can change which one you are attuned to during a long rest.

If you lose your Stormforged Weapon or wish to create additional ones, you can do so over the course of three days (8 hours each day) by expending 200 gp worth of metal and other raw materials. When you make a new Stormforged Weapon, you can make the same or a different type, and select the same or different upgrades.

Select one of the following and consult the Stormforged Weapon table for its statistics.
\savebox{\StormforgedBox}{%
  \begin{varwidth}{\dimexpr\textwidth + 1cm}
    \small
    \renewcommand{\arraystretch}{1.15}
    \rowcolors{2}{AltRow}{Parchment}
    \begin{tabularx}{\hsize}{
	  @{}
	  >{\raggedright\arraybackslash\hyphenpenalty=10000\exhyphenpenalty=10000\relax}p{2.7cm}
	  >{\raggedright\arraybackslash\hyphenpenalty=10000\exhyphenpenalty=10000\relax}p{2.1cm}
	  >{\centering\arraybackslash}p{1.4cm}
	  >{\raggedright\arraybackslash}p{1.7cm}
	  >{\raggedright\arraybackslash}X
	  @{}
	}

    \rowcolor{TableHeader}
    \textbf{Weapon Name} & \textbf{Damage} & \textbf{Weight} & \textbf{Type} & \textbf{Properties} \\
    Thunder Cannon & 1d12 piercing & 15 lbs. & Ranged & Ammunition (60/180), Two-Handed, Loud², Stormcharged¹ \\
    Hand Cannon & 1d10 piercing & 5 lbs. & Ranged & Ammunition (30/90), Light, Loud², Stormcharged¹ \\
    Kinetic Hammer & 1d10 bludgeoning + 1d4 thunder & 10 lbs. & Melee & Two-Handed, Heavy, Loud² \\
    Charged Blade & 1d6 slashing + 1d4 lightning & 3 lbs. & Melee & Finesse, Loud² \\
    Lightning Pike & 1d8 piercing + 1d4 lightning & 10 lbs. & Melee & Reach, Two-Handed, Loud² \\
    \end{tabularx}
  \end{varwidth}
}

\begin{center}
\begin{tikzpicture}
  \node[
    fill=Parchment, 
    draw=black, 
    rounded corners=2pt, 
    inner sep=6pt,
    anchor=north west,
  ] (box) at (0,0) {\usebox{\StormforgedBox}};
\end{tikzpicture}
\end{center}


\vspace{2pt}
\footnotesize
\noindent
\textsuperscript{1}\textbf{Stormcharged.} When you use an action, bonus action, or reaction to attack with a Stormcharged Weapon, you can make only one Attack regardless of the number of attacks you can normally make. If you could otherwise make additional attacks with that action, the weapon deals an extra 3d6 lightning or thunder damage per attack that was foregone. \par
\textsuperscript{2}\textbf{Loud.} Your weapon rings with thunder that is audible within 300 feet of you whenever it makes an attack.




\Subheading{Thunder Cannon}

You use the power of thunder to launch a projectile with terrible power, if limited accuracy, over long distances. Ringing out with a booming crash, it brings fear to the battlefield.
\vspace{1in}

\Subheading{Hand Cannon}

Forgoing the guiding barrel, this pack uses the thundering power to launch a projectile with all the force of a Cannon, though its effective range is far more limited.

\Subheading{Kinetic Hammer}

Rather than launching a projectile with the thundering force, you keep that force imbued in the weapon, allowing for devastating force to be applied to the attack.

\begin{tikzpicture}[remember picture, overlay]
  \node[
    anchor=south east,
    inner sep=0pt
  ] at ([xshift=1.1cm, yshift=1.5cm]current page.south east)
  {\includegraphics[width=9.0cm]{Stormcharged.png}};
\end{tikzpicture}
\newpage

\Subheading{Charged Blade}

You create a bladed weapon that channels the harnessed power of the elemental storm power directly into the blade, causing it to lay waste to all it strikes. This weapon deals lightning damage instead of thunder when applying Thundermonger.

\begin{CalloutBox}{Firearms in a Campaign Setting}

A lot of campaign settings do not feature firearms, and in some of these the Thundersmith variant of inventor might not be the right choice, but consider that the inventor is fundamentally someone that tinkers with and explores boundaries of magic as much as or more so than technology.

A Thunder Cannon need not be a gunpowder powered device, even in a setting where gunpowder exists, but can be powered by harnessing elemental powers, bound through various carefully researched magical techniques. In most cases, the wonders of an inventor are more an engineering marvel of magic than technology, but that balance can shift depending on what is best for your setting.
\end{CalloutBox}

\Subheading{Lightning Pike}

You create a charged blade and stick it to the end of a pole, making it slightly more unwieldy, but giving it devastating reach. This weapon deals lightning damage instead of thunder when applying Thundermonger.

\subsection{Thundermonger}

At 3rd level, the elemental might of your weapon is so powerful its strikes deal bonus thunder damage.

When you hit a target with your Stormforged Weapon, you can deal an extra 1d6 thunder damage. After discharging this bonus damage, you can’t deal this bonus damage again until the start of your next turn.

This extra damage increases by 1d6 when you reach certain levels in this class: 5th level (2d6), 7th level (3d6), 9th level (4d6), 11th level (5d6), 13th level (6d6), 15th level (7d6), 17th level (8d6), and 19th level (9d6).

\begin{CalloutBox}{Stormcharged vs. Thundermonger}

A common question is why does the Stormcharged Property exist and how does it interact with Thundermonger? Functionally, the stormcharged property (like the loading property) has no effect on a single classed Thundersmith, it just exists to make Stormforged Weapons interact better with Extra Attack for multiclassing and Cross-Disciplinary Knowledge.
\end{CalloutBox}

\subsection{Devastating Blasts}

Beginning at 5th level, when you miss an attack with your Thundering Weapon, you can apply Thundermonger damage to the target creature you missed, but it deals only half the bonus damage. Dealing damage this way counts as applying Thundermonger damage.

\subsection{Unleashed Power}

Starting at 14th level, when rolling damage for Thundermonger or your Stormforged Weapon, you can expend a spell slot to reroll a number of the damage dice up to your Intelligence modifier (minimum of one) and maximize a number of dice equal to the level of the spell slot expended. (You may maximize the dice rerolled this way, or the dice that were not rerolled) You must use the new rolls.

\section{Thundersmith Upgrades}

\subsection{Unrestriced Upgrades}

\begin{CalloutBox}{How to Merge Your Weapons}

If you want to call it a gunblade, I’m not here to stop you, but it can be any range of configuration, from a bayonet to something more exotic.
\end{CalloutBox}

\Subheading{Adaptable Weapon }

You can adapt a weapon without the ammunition property to have an alternate attack type giving it the functionality of a Hand Cannon, or give a weapon with the ammunition property an alternate attack with the properties of a Charged Blade.

\Subheading{Extended Range}
 
You extend the reach of your Stormforged Weapon. If your Stormforged Weapon has the ammunition property, its range is extended; normal range is extended by 30 feet, and maximum range is extended by 90 feet. If your Stormforged Weapon is a melee weapon with the two-handed property, it gains the reach property.

\Subheading{Lightning Burst }

You upgrade your Stormforged Weapon to discharge its power within a 5-foot-wide 60-foot-long line. If you have not dealt Thundermonger damage since the start of your turn, as an action, you can make a special attack. Each creature in the area must make a Dexterity saving throw against your spell save DC or take damage equal to the bonus damage of Thundermonger as lightning damage on a failed save, half as much on a successful save.

This counts as discharging your Thundermonger damage. Firing in this way doesn’t consume ammo.

\Subheading{Lightning Magic}
 
After a long study of the internal workings of your Stormforged Weapon, your mastery of lightning and thunder magic is such that you learn the following spells at the following levels and can cast them as Inventor Spells.

\vspace{1pt}
\begin{SmallTableBox}
\small % optional, helps in two-column mode
\renewcommand{\arraystretch}{1.2}
\rowcolors{2}{AltRow}{Parchment}
\begin{tabularx}{\linewidth}{@{}>{\raggedright\arraybackslash}p{2.5cm}X@{}}
\rowcolor{TableHeader}
\textbf{Inventor Level} & \textbf{Spell} \\
3rd  & \textit{lightning tendril*} \\
5th  & \textit{crackle*} \\
9th  & \textit{lightning bolt} \\
13th & \textit{jumping jolt*} \\
15th & \textit{sky burst*} \\
\end{tabularx}
\end{SmallTableBox}
\vspace{1pt}

You can cast lightning tendril* without expending a spell slot, after which you must finish a long rest before you can cast it without expending a spell slot again. Starting at 5th level, you can choose to cast crackle* without expending a spell slot instead.

\Subheading{Point Blank (Prerequisite: Hand Cannon)}
 
Being within 5 feet of a hostile creature doesn’t impose disadvantage on your ranged attack rolls. Additionally, you can use your Hand Cannon when making opportunity attacks.

\Subheading{Silencer (Incompatible with Echoing Boom)}
 
You upgrade your Stormforged Weapon with a sound dampening module. Your Stormforged Weapon loses the loud property. Additionally, you can expend a 2nd-level spell slot to overcharge the Silencer, casting the silence spell

\textbf{Note:}  \textit{Sound—and consequently thunder damage—can’t pass through silence or affect a creature inside its area of effect.}

\Subheading{Shock Absorber}
 
You add a reclamation device to your Stormforged Weapon to gather energy from the surroundings when it is present. As a reaction to taking Lightning or Thunder damage, you can gain resistance to that damage type until the start of your next turn. When you do so, the next time you roll Thundermonger damage before the end of your next turn, it deals an extra 1d6 damage.

\Subheading{Sonic Movement (Prerequisite: Stormforged weapon that can deal Thunder damage)}
 
You recalibrate and rebalance your weapon to leverage the backdraft of the force it exerts. When you make an attack with your Stormforged weapon that would deal thunder damage on your turn, you are knocked 5 feet away from the target you attacked. This movement doesn’t provoke opportunity attacks.

\Subheading{Storm Blast}
 
You can make a special attack with your Stormforged Weapon, unleashing a storm blast in a 30-foot cone as an action. Each creature in the area must make a Strength saving throw against your spell save DC, or take damage equal to the bonus damage of Thundermonger and be knocked prone.

This counts as discharging your Thundermonger damage. Firing in this way doesn’t consume ammo.

If using a Kinetic Hammer, you can use your Strength to calculate the DC of this ability (8 + your Strength modifier + your proficiency bonus) instead of your spell save DC.

\Subheading{Twin Thunder}
 
You can attune to two one-handed Stormforged Weapons at the same time, so long as they either share the same Upgrades, or have total Upgrades equal to your maximum Upgrade count between them. Attuning to a second weapon this way doesn’t count against your maximum number of attuned items. If you make an attack with one of them while holding the other, you can attack with the other as a bonus action, however, both share the same use of Thundermonger. You don’t add your ability modifier to the damage of the bonus attack, unless that modifier is negative or you have the Two Weapon Fighting Fighting Style.

While dual wielding Stormforged Weapon, you can load a Stormforged Weapon without a free hand.

\Subheading{Weapon Improvement}
 
Your Stormforged Weapon gains a +1 to attack and damage rolls. This doesn’t stack with any benefit gained from Arcane Retrofit, and this upgrade can be replaced as part of applying a bonus to your Stormforged Weapon via Arcane Retrofit.

\subsection{5th-Level Upgrades}

\Subheading{Echoing Boom (Incompatible with Silencer)}
 
You pack extra power into your Thundermonger, increasing the damage it deals by 1d6.

\Subheading{Harpoon Reel}
 
You devise an alternate attack method for your Stormforged Weapon that launches a harpoon attached to a tightly coiled cord. This attack has a normal range of 30 feet and a maximum range of 60 feet, and it deals 1d6 piercing damage. This attack doesn’t apply Thundermonger damage. This attack can target a surface, object, or creature.

A creature struck by this attack is impaled by the Harpoon unless it removes the Harpoon as an action, which causes it to take an extra 1d6 damage. While the Harpoon is stuck in the target, you are connected to the target by a 60 foot cord. Dragging the connected party via the attached cord causes the creature moving to move at half speed unless they are a size category larger.

While connected in this manner, you can use your bonus action to activate the Reel action, pulling yourself to the location if the target is Medium or larger. A Small or smaller creature is pulled back to you, ending the connection. Alternatively, you can opt to disconnect the cord.

This attack can’t be used again until the Reel action is taken.


\begin{CalloutBox}{Connected...By a Rope}

Note that Harpooning a creature means that there is a rope connecting you to them. While you can use this to move them as per the ability, they can also use this to move you, particularly if they are substantially larger. Many Thundersmiths were last seen saying “I got this” and then Harpooning a dragon. User discretion is advised.
\end{CalloutBox}

\Subheading{Terrifying Thunder (Prerequisite: Echoing Boom)}
 
You add an additional amplifier to maximize the shock and awe value of your Stormforged Weapon. The first time a target takes damage from Thundermonger, they are deafened until the end of their next turn.
Additionally, they must make a Wisdom saving throw against your spell save DC or become frightened of you for one minute. They can repeat this saving throw at the end of each of their turns.

On a successful saving throw, the creature is immune to Terrifying Thunder for 24 hours.

\subsection{9th-Level Upgrades}

\Subheading{Elemental Swapping}
 
You upgrade the firing chamber for more adaptable damage. When you take the attack action with your Stormforged Weapon you can adjust its elemental properties, causing any bonus damage granted by Thundermonger to deal acid, cold, fire, lightning, or thunder damage instead of thunder damage.

Alternatively, you can use a vial of holy water to cause your next Thundermonger bonus damage to deal radiant damage.

\Subheading{Mortar Shells (Prerequisite: Stormforged weapon with the Ammunition property)}
 
You build an alternate fire mode allowing you to fire your Stormforged Weapon like a mortar. Pick a target point within range, and make an attack roll. Apply the attack roll to all creatures within a 5-foot radius of the target point. Creatures hit take the weapon's damage plus half of the Thundermonger bonus damage. Dealing damage this way counts as applying Thundermonger damage for the turn. If you miss all targets, you can apply Devestating Blasts as normal.

Creatures do not benefit from cover against this fire mode unless they have overhead cover as well.

\Subheading{Ride the Lightning (Prerequisite: Lightning Burst)}
 
When you use the Lightning Burst ability, you can opt to expend a spell slot of 1st level or higher to infuse yourself into the burst of power. When you do so you are teleported up to 60 feet in the direction of the Lightning Burst (ending early if the Lighting Burst is blocked by an obstacle). You can stop anywhere along the path of the Lighting Burst, but the line will stop where you do.

\Subheading{Shock Harpoon (Prerequisite: Harpoon Reel)}
 
After hitting a creature with the Harpoon fire mode, you can use a bonus action to deliver a shock. If you have not dealt Thundermonger damage since the start of your turn, you can deal damage equal to your Thundermonger bonus damage as lightning damage. This counts as discharging your Thundermonger damage. Additionally, the target must make a Constitution saving throw against your spell save DC or be stunned until the end of its next turn.

Once used, the Harpoon must be reeled in before this can be used again.

\Subheading{Synaptic Feedback}
 
You install a feedback loop into your cannon, allowing you to siphon some energy from your Stormforged Weapon to empower your reflexes.
Whenever you deal lightning damage with your Stormforged Weapon your walking speed increases by 10ft and you can take the Dash or Disengage actions as a bonus action. This effect lasts until the end of your turn.

\Subheading{Thunder Jump}
 
You build a quick release for the arcane thundering energy that fills your Stormforged Weapon. As an action you can discharge your Thunder Cannon downward, launching yourself up to 60 feet in any direction. This movement doesn’t provoke opportunity attacks, and you carry up to one willing creature of your choice within 5 feet with you. Neither you or the carried creature takes falling damage on landing.

Creatures that are within 10 feet of you when you use this ability must make a Dexterity saving throw against your spell save DC. On failure, they take damage equal to half of your Thundermonger damage. On success, they take half as much damage. This counts as discharging your Thundermonger damage. Firing in this way doesn’t consume ammo.

\Subheading{Transforming Weapon (Prerequisite: Adaptable Weapon)}
 
Select 3 of Thunder Cannon, Hand Cannon, Kinetic Hammer, Charged Blade, and Lightning Pike. As a bonus action, you can transform your Stormforged Weapon between them.

\subsection{11th-Level Upgrades}

\Subheading{Backblast}
 
After dealing thunder damage on your turn, you can create a burst of thunderous sound that can be heard up to 100 feet away as a bonus action. Each creature within 5 feet, other than you, must make a Constitution saving throw or take 3d6 thunder damage You can do this a number of times equal to your Intelligence modifier, regaining all uses after a long rest. This damage increases by 1d6 when you reach 17th level to a total of 4d6.

\Subheading{Blast Radius}
 
Your Devastating Blasts feature now deals half your weapon damage (including your modifier) in addition to half your Thundermonger damage when you apply it to a missed target that is within 30 feet of you.

\Subheading{Stabilization}
 
When making a ranged weapon attack with your Stormforged Weapon, if neither you nor your target has moved since your last ranged attack against them, then you gain advantage on the attack. Also, being prone no longer causes you to have disadvantage when making a ranged weapon attack with your Stormforged Weapon.

\subsection{15th-Level Upgrades}

\Subheading{Massive Overload (Prerequisite: Storm Blast or Lightning Burst)}
 
Before making an attack with your Stormforged Weapon, you can expend a spell slot of 3rd level or higher to use Storm Blast or Lightning Burst at the same time as making an attack; in this case the Storm Blast or Lightning Burst is powered by the spell slot and doesn’t count as applying Thundermonger for that turn. The direction of this secondary ability is the same as your attack.

Doing this damages your Stormforged Weapon and you must spend an action to repair it before you can fire again. You must have the secondary ability unlocked as an upgrade to use it.

\Subheading{Masterwork Weapon (Prerequisite: Weapon Improvement)}
 
The bonus to attack and damage rolls for your Stormforged Weapon increases by +2 (stacking with any existing bonus from Weapon Improvement or Arcane Retrofit), up to a maximum of +4.

\SectionSplashPage[0.4\textheight]{Warsmith.png}
\begin{multicols}{2}
\section{Warsmith}
A Warsmith is an Inventor that has turned their wondrous talent of invention to a singular goal: making themselves a juggernaut of war. The reasons behind this could be benevolent or nefarious. Some Warsmiths seek to turn their invention into a machine of death and terror; others become the arbiter of justice and order, and others still perhaps merely seek to refine their craft in pursuit of pure innovation.

Because few individuals would pursue such a wondrous invention without a driven purpose to their endeavor, warsmiths tend to be lawful. Usually driven to their actions by a greater purpose, they seek the power to accomplish their aims, be it righting the wrongs of the world or bringing it to heel beneath their ironshod boot.

\subsection{Warsmith's Proficiency}

At 1st level, you gain proficiency with heavy armor, tinker’s tools and smith’s tools.

\subsection{Warplate Gauntlet}

At 1st level, when you take this specialization, you construct a Warplate Gauntlet. This is a specialized Wondrous Item that only you can attune to. When you create a Warplate Gauntlet, you can add one of the following upgrades to it: Power Fist, Force Blast, or Martial Grip. This upgrade doesn’t count against your upgrade total. You can make multiple gauntlets with different upgrades, but can only be attuned to one at a time.

If you lose your Warplate Gauntlet, you can remake it during a long rest with 25 gp worth of materials, or can scavenge for materials and forge it over two days of work (8 hours each day) without the material expense.

While wearing a Warplate Gauntlet, you can engage Artificial Strength.
\end{multicols}
\newpage
\twocolumn

\Subheading{Artificial Strength}

When you don your Warplate Gauntlet or as an action while wearing it, you can dedicate some of your intelligence to fully controlling the power of the gauntlet. You can reduce your current and maximum Intelligence score to increase your current Strength ability score by the same amount, but you can only raise your Strength ability score up to what your Intelligence ability score was before engaging Artificial Strength. You can stop using Artificial Strength at any time, and it automatically ends if your gauntlet is removed.

\subsection{Warsmith's Armor}

At 3rd level, you’ve attained the knowledge of forging and arcane tinkering sufficient to create a set of armor that augments and expands your abilities from a standard, nonmagical, set of heavy armor using resources you’ve gathered. This process takes 8 hours to complete, requiring the use of a forge, foundry or similar and it incorporates a Warplate Gauntlet (they do not require separate attunement). While wearing your armor, your Strength ability score increases by 2, and your maximum Strength ability score becomes 22.

You can create a new set of armor by forging it from a set of gathered and purchased materials in a process that takes 2,000 gp and 8 hours. You can create multiple sets of armor, but you can only be attuned to one of them at a given time, and you can only change which one you are attuned to during a long rest. If you create a new set of Warsmith’s Armor, you can apply a number of Upgrades equal to the value on the class table, applying each at the level you get it on the class table.

When you create your armor, you can create a heavy plated Warplate, medium balanced Wargear, light flexible Warsuits, or you can integrate your changes directly into your body as Integrated Armor. If your armor is Warplate or Integrated Armor, you gain the Powerful Build trait. Powerful Build means you count as one size larger when determining your carrying capacity and the weight you can push, drag, or left. Integrated Armor can’t be removed, but doesn’t confer any penalty when resting in it.

Additionally, if you are Small, you become Medium while wearing Warplate.

\Subheading{Warsmiths \& Magical Armor}

By the rules laid out here, using magical armor as a base for your Warsmith’s armor has no additional effect. This is intentionally the rules-as-written, but there is certainly some flexibility here.

\begin{fiveitemize}
  \item Using Adamantine or Mithral, the properties carry over to the Warsmith’s armor.
  \item Using +1/+2/+3 armor carries over, but counts as taking a free “Armor Class” upgrade for each +1 the armor has, meaning that upgrade can’t be taken to make the armor +4 or better.
  \item Armor of Resistance carries over, counting as taking “Resistance” upgrade for that damage type for free.
\end{fiveitemize}

Other cases can be handled on a case-by-case basis. Consult with your GM and work something out that would be reasonable to combine making receiving magic armor a cool bonus, but not something that breaks the game!

\savebox{\WarsmithBox}{%
  \begin{varwidth}{\dimexpr\textwidth + 1cm}
    \small
    \renewcommand{\arraystretch}{1.15}
    \rowcolors{2}{AltRow}{Parchment}
    \begin{tabularx}{\hsize}{
	  @{}
	  >{\raggedright\arraybackslash}p{2.7cm}
	  >{\raggedright\arraybackslash\hyphenpenalty=10000\exhyphenpenalty=10000\relax}p{2.1cm}
	  >{\centering\arraybackslash}p{1.4cm}
	  >{\raggedright\arraybackslash}p{1.7cm}
	  >{\raggedright\arraybackslash}X
	  @{}
	}

    \rowcolor{TableHeader}
    \textbf{Armor Name} & \textbf{AC} & \textbf{Weight} & \textbf{Stealth} & \textbf{Properties} \\
    Warplate & 18 & 65 lbs. & Disadvantage & Heavy Armor, Powerful Build¹ \\
    Warsuit & 14 + Dex Modifier (max 2) & 20 lbs. & — & Medium Armor \\
    Warskin & 12 + Dex Modifier & 13 lbs. & — & Light Armor \\
    Integrated Armor & 14 + Dex modifier (max 2) & N/A & — & Medium Armor, Powerful Build¹ \\
    \end{tabularx}
  \end{varwidth}
}

\begin{center}
\begin{tikzpicture}
  \node[
    fill=Parchment, 
    draw=black, 
    rounded corners=2pt, 
    inner sep=6pt,
    anchor=north west,
  ] (box) at (0,0) {\usebox{\WarsmithBox}};
\end{tikzpicture}
\end{center}


\vspace{2pt}
\footnotesize
\noindent
\textsuperscript{1}\textbf{Powerful Build.} You count as one size larger when determining your carrying capacity and the weight you can push, drag, or lift. \par
\end{document}